\documentclass[sigconf,anonymous]{acmart}

\setcopyright{none}
\copyrightyear{2026}
\acmYear{2026}
\acmDOI{}
\acmISBN{}
\acmConference[Conference '26]{Conference}{June 2026}{City, Country}

\usepackage{amsmath,amssymb,amsthm}
\usepackage{algorithm}
\usepackage{algpseudocode}
\usepackage{booktabs}
\usepackage{listings}
\usepackage{xcolor}
\usepackage{enumitem}
\usepackage{subcaption}
\usepackage{multirow}
\usepackage{cite}
\usepackage{url}
\usepackage{graphicx}
\usepackage{balance}
\usepackage{pifont}

% ── Listing style ──────────────────────────────────────────
\lstset{
  basicstyle=\ttfamily\footnotesize,
  keywordstyle=\bfseries,
  commentstyle=\itshape\color{gray},
  breaklines=true,
  frame=single,
  numbers=left,
  numberstyle=\tiny\color{gray},
  xleftmargin=1.6em,
  framexleftmargin=1.4em,
  columns=flexible,
  captionpos=b,
  aboveskip=6pt,
  belowskip=4pt
}

% ── Theorem environments ──────────────────────────────────
\newtheorem{definition}{Definition}
\newtheorem{theorem}{Theorem}
\newtheorem{proposition}{Proposition}
\newtheorem{lemma}{Lemma}

\begin{document}

\title{CIR\,+\,CVN: Bridging LLM Semantic Understanding\\
and Petri-Net Verification for Concurrent Programs}

\author{
\IEEEauthorblockN{[Author Names]}
\IEEEauthorblockA{[Institution]\\[City, Country]\\[email]}
}


% ══════════════════════════════════════════════════════════════
%  ABSTRACT
% ══════════════════════════════════════════════════════════════
\begin{abstract}
% 通过静态分析检测并发错误从根本上受限于别名问题：确定两个引用是否指向同一个共享资源通常是不可判定的，然而精确的别名信息却是锁序分析、数据竞争检测和同步协议验证的先决条件。我们发现，通过重新构建问题可以完全绕过这一障碍。我们不再向静态分析器询问"哪些程序引用存在别名？", 而是向逻辑逻辑模型（LLM）询问"该程序打算共享哪些资源？"——这是一个关于业务语义的问题，LLM 可以可靠地回答这个问题。基于此，我们提出了一种包含两种新形式化的"生成-验证-修复"架构。 (1) 并发中间表示 (Cir)：一种由 LLM 生成的语句级中间语言，其中每个共享资源都被全局命名，并且每个锁-变量保护关系都被声明，从而通过构造而非分析消除别名。(2) 并发验证网 (Cvn)：一种特定领域的加权 P/T Petri 网，具有全局变量保护和三值求值，它仅使用单色黑色标记对所有同步相关行为进行建模。Cir 被机械地转换为 Cvn，后者通过穷举状态空间搜索来发现死锁、信号丢失、饥饿和活锁。结构化的反例被映射回 Cir 语句标识符，并作为目标修复提示提供给 LLM，从而完成"生成-验证-修复"循环。我们形式化了这两种表示，证明了 \textsc{Cir} 执行和 \textsc{Cvn} 触发之间的步骤对应定理，定义了一个具有 61~静态检查规则和 5~错误类型检测器的两层验证架构，并在涵盖死锁、信号丢失和饥饿的 10~代表性并发模式上验证了该流水线。
Detecting concurrency bugs through static analysis is fundamentally limited by the \emph{alias problem}: determining whether two references denote the same shared resource is undecidable in general, yet precise alias information is a prerequisite for lock-order analysis, data-race detection, and synchronization-protocol verification. We observe that this barrier can be circumvented entirely by reframing the question. Instead of asking a static analyzer ``which program references happen to alias?'', we ask an LLM``which resources does this program intend to share?''---a question about business semantics that LLMs answer reliably. Building on this observation, we propose a \emph{generate--verify--repair} architecture comprising two new formalisms. (1)~\emph{Concurrency Intermediate Representation} (\textsc{Cir}): an LLM-generated, statement-level intermediate language in which every shared resource is globally named and every lock--variable protection relation is declared, eliminating aliases by construction rather than by analysis. (2)~\emph{Concurrency Verification Net} (\textsc{Cvn}): a domain-specific weighted P/T Petri net with global-variable guards and three-valued evaluation, which models all synchronization-relevant behavior using only single-color black tokens. \textsc{Cir} is mechanically translated to \textsc{Cvn}, whose exhaustive state-space search discovers deadlocks, signal losses, starvation, and live lock. Structured counterexamples are mapped back to \textsc{Cir} statement identifiers and fed to the LLM as targeted repair prompts, closing the \emph{generate--verify--repair} loop. We formalize both representations, prove a step-correspondence theorem between \textsc{Cir} executions and \textsc{Cvn} firings, define a two-layer verification architecture with 61~static-check rules and 5~bug-kind detectors, and validate the pipeline on 10~representative concurrency patterns covering deadlocks, signal losses, and starvation.
\end{abstract}

\keywords{concurrency verification, Petri nets, large language models,
alias analysis, intermediate representation, deadlock detection}

\maketitle

% ══════════════════════════════════════════════════════════════
%  Section files
% ══════════════════════════════════════════════════════════════
% ══════════════════════════════════════════════════════════════
%  INTRODUCTION
% ══════════════════════════════════════════════════════════════
\section{Introduction}
\label{sec:introduction}

% ── 1. 并发 bug 的严重性 ──────────────────────────────────
% 并发错误——包括死锁、数据竞争、信号丢失和饥饿——是软件系统中代价最高的缺陷之一。它们源于线程、锁、条件变量、通道和共享状态之间微妙的交互，并且由于依赖于调度的不确定性，因此极难复现。与确定性地显现的顺序错误不同，并发错误可能在数百万次的测试运行中都难以被发现，而只有在特定的时间条件下才会在生产环境中出现。工业界和学术界都采用了两种主要的策略来对抗这些错误：动态分析，即在运行时探索程序的执行过程；以及静态分析，即在不执行程序的情况下推断其行为。诸如 Miri、ThreadSanitizer 和 Loom 之类的动态工具能够有效地检测已探索调度中的错误，但它们本质上无法保证未探索调度中不存在错误。原则上，静态分析可以提供这种保证——但它面临着一个根本性的障碍。
Concurrency bugs---deadlocks, data races, signal losses, and starvation---are among the most costly defects in software systems. They arise from subtle interactions among threads, locks, condition variables, channels, and shared state,and they are notoriously difficult to reproduce because they depend on scheduling nondeterminism. Unlike sequential bugs that manifest deterministically, a concurrency bug may hide through millions of test runs and surface only in production under specific timing conditions.Both industry and academia have pursued two broad strategies to combat these bugs: \emph{dynamic analysis}, which explores program executions at runtime, and \emph{static analysis}, which reasons about program behavior without executing it. Dynamic tools such as Miri, Thread Sanitizer, and Loom are effective at detecting bugs along explored schedules, but they inherently cannot guarantee the absence of bugs along unexplored schedules. Static analysis, in principle, can provide such guarantees---but it faces a fundamental barrier.

% ── 2. 别名问题：静态分析的根本障碍 ─────────────────────
% 根本障碍在于别名问题：确定两个程序引用是否指向同一内存位置。这个问题通常是不可判定的，即使对于顺序程序，在实践中也难以精确判断。对于并发程序，情况则更为糟糕。考虑一个 Rust 程序，其中两个线程各自接收一个 Arc<(Mutex<State>, Condvar)> 的克隆。为了确定两个线程是否操作同一个互斥锁-条件变量对，静态分析器必须追踪该值经过堆分配、Arc::clone、模式解构、借用和闭包捕获的过程——这条链使得大多数指针分析都无法奏效。然而，别名信息对于并发验证来说并非可有可无。如果不知道两个锁操作是否指向同一个互斥锁，死锁检测器就无法构建锁序图。如果不知道变量访问和锁获取指向一个受保护的变量对，竞争检测器就无法判断该访问是否受保护。所有主要的并发分析类别——锁顺序分析、先行断言、数据竞争检测、同步协议验证——都要求精确的别名信息作为前提条件。因此，别名问题不仅仅是造成精度不足的原因；它是整个分析的“承重墙”，一旦移除，整个分析就会崩溃。现有工具通过各种折衷方案来应对这一障碍。流不敏感分析通过合并所有程序点的别名集来牺牲精度。流敏感分析可以实现更高的精度，但扩展性差。动态工具通过在运行时观察具体地址来绕过别名，但会牺牲覆盖率。所有这些折衷方案都无法满足“自动化修复”工作流程的要求，因为在自动化修复工作流程中，每个诊断结果都必须既精确（避免浪费修复工作的误报）又可操作（指向具体的代码位置）。
\subsection{The Alias Barrier in Concurrency Analysis}
\label{subsec:alias-barrier}
The fundamental barrier is the \emph{alias problem}: determining whether two program references denote the same memory location. This problem is undecidable in general and remains imprecise in practice even for sequential programs. For concurrent programs, the situation is substantially worse. Consider a Rust program in which two threads each receive a clone of an\texttt{Arc<(Mutex<State>,\,Condvar)>}. To determine whether both threads operate on the \emph{same} mutex--condvar pair, astatic analyzer must trace the value through heap allocation,\texttt{Arc::clone}, pattern destructuring, borrowing, and closure capture---a chain that defeats most pointer analyses.Yet alias information is not optional for concurrency verification.Without knowing that two lock operations target the same mutex, a deadlock detector cannot construct a lock-order graph. Without knowing that a variable access and a lock acquisition refer to a guarded pair, a race detector cannot determine whether the access is protected. Every major category of concurrency analysis---lock-order analysis, happens-before reasoning, data-race detection,synchronization-protocol verification---requires precise alias information as a prerequisite. The alias problem is therefore not merely a source of imprecision; it is the \emph{load-bearing wall}whose removal collapses the entire analysis.Existing tools cope with this barrier through various compromises.Flow-insensitive analyses sacrifice precision by merging alias sets across all program points. Flow-sensitive analyses achieve better precision but scale poorly. Dynamic tools sidestep aliases by observing concrete addresses at runtime, but sacrifice coverage. None of these compromises is satisfactory for an\emph{automated repair} workflow in which every diagnostic must be both precise (no false positives that waste repair effort) and actionable (pointing to specific code locations).

% ── 3. 关键洞察：不是解析别名，而是理解业务 ──────────────
% 我们观察到，别名问题的难点在于其框架。传统分析提出的问题是关于实现工件的：“给定这堆指针、闭包和引用计数句柄，哪些表达式对恰好指向同一个对象？”这是一个关于偶然巧合的问题，其难点恰恰在于答案取决于指针操作的完整历史。但对于并发分析而言，有一个更简单的问题也能提供相同的信息：“这个程序打算使用哪些共享资源，哪些锁保护哪些变量？”这是一个关于业务语义的问题——关于程序员打算构建什么，而不是关于编译器看到什么。
\subsection{Reframing: From ``What Aliases'' to ``What Should Be Shared''}
\label{subsec:reframing}
We observe that the difficulty of the alias problem stems from its framing. Traditional analysis asks a question about \emph{implementation artifacts}: ``Given this heap of pointers, closures, and reference-counted handles, which pairs of expressions happen to refer to the same object?'' This is a question about emergent coincidence, and it is hard precisely because the answer depends on the full history of pointer manipulation. But there is a much easier question that yields the same information for concurrency analysis: ``What shared resources does this program \emph{intend} to use, and which locks protect which variables?'' This is a question about \emph{business semantics}---about what the programmer \emph{meant} to build, not about what the compiler sees.

% 基于海量代码语料库训练的大型语言模型（LLM）擅长此类语义理解。读取生产者-消费者实现的LLM无需追踪Arc::clone链；它理解程序有一个共享队列、一个保护它的互斥锁和一个指示数据可用性的条件变量。它之所以知道这一点，是因为它理解了业务模式，而不是因为它执行了指针分析。这一洞见提示了一种根本性的重新构建：与其构建一个更好的别名分析器，不如让LLM生成一个不存在别名的表示。如果每个共享资源在生成时都获得一个唯一的全局名称，则无需解析别名——问题通过构造而非分析得到消除。LLM的作用不是“解析别名”（这相当于要求它执行与传统分析相同的艰巨任务）；它的作用是根据对业务逻辑的理解来声明程序的资源架构。
Large language models (LLMs), trained on massive code corpora, excel at precisely this kind of semantic understanding. An LLM reading a producer--consumer implementation does not need to trace \texttt{Arc::clone} chains; it understands that the program has one shared queue, one mutex protecting it, and one condvar signaling data availability. It knows this because it understands the \emph{business pattern}, not because it performs pointer analysis. This insight suggests a fundamental reframing: instead of building a better alias analyzer, \textbf{we ask the LLM to generate a representation in which aliases do not exist}. If every shared resource receives a unique global name at generation time, there are no aliases to resolve. The LLM's role is not to resolve aliases which would be asking it to perform the same hard task as traditional analysis; its role is to declare the resource architecture of the program based on its understanding of the business logic.

% ── 4. CIR：基于业务理解的中间表示 ───────────────────────
% 基于这种重新构建，我们提出了并发中间表示（Cir）。Cir 是一种中间语言，旨在由逻辑逻辑模型（LLM）生成，并由形式验证引擎使用。其核心设计原则直接解决了别名障碍：全局资源命名。每个共享资源——互斥锁、读写锁、条件变量、信号量、通道、共享变量、原子变量——都拥有一个唯一的全局名称。函数直接通过名称引用资源，无需参数或返回值。没有指针、堆，也没有间接寻址。显式保护映射。Cir 在一个专门的“保护”部分声明哪些锁保护哪些变量。这种关系在传统分析中需要费力推断，而在这里则被声明为一等声明，因为 LLM 可以通过理解程序设计自动获知它。语句级标识符。每个 Cir 语句都带有一个唯一的标识符 (sid)，为反例报告和定向修复提供稳定的锚点。Cir 不是通用编译器中间表示 (IR)。它仅保留与并发验证相关的结构：同步操作、线程控制和共享状态上的分支。业务逻辑（算术运算、字符串操作、I/O）已被抽象化。这种刻意的缩小范围使得 LLM 生成可靠：LLM 不需要生成完整的可执行程序；它只需要捕获并发架构。静态检查器会根据 61 条规则（分为 9 个错误类别 E0xx 到 E8xx）验证 Cir 工件，在执行更耗时的验证步骤之前捕获结构错误（缺少字段、未定义的资源）、语义错误（对非锁定资源加锁、未受保护的变量访问）和一致性错误（spawn/join 不匹配、不可达语句）。
\subsection{\textsc{Cir}: Concurrency Intermediate Representation}
\label{subsec:intro-cir}

Based on this reframing, we propose the \emph{Concurrency Intermediate Representation} (\textsc{Cir}). \textsc{Cir} is an intermediate language designed to be generated by an LLM and consumed by a formal verification engine. Its core design principles directly address the alias barrier: 1. Every shared resource (mutex, read-write lock, condition variable, semaphore, channel, shared variable, atomic variable) receives a unique global name. Functions reference resources directly by name without parameters or return values. There are no pointers, no heap, no indirection. 2. The \textsc{Cir} declares which locks protect which variables in a dedicated \texttt{protection} section. This relationship, which traditional analysis must laboriously infer, is stated as a first-class declaration because the LLM knows it from understanding the program's design. 3. Each \textsc{Cir} statement carries a unique identifier ($\mathsf{sid}$), providing stable anchors for counterexample reporting and targeted repair. 

\textsc{Cir} is not a general-purpose compiler IR. It retains only the structure relevant to concurrency verification: synchronization operations, thread control, and branching over shared state. Business logic (arithmetic, string manipulation, I/O) is abstracted away. This deliberate narrowing is what makes LLM generation reliable: the LLM does not need to produce a complete, executable program; it only needs to capture the concurrency architecture. A static checker validates \textsc{Cir} artifacts against 61~rules organized in 9~error categories, catching structural errors (missing fields, undefined resources), semantic errors (lock on non-lock resource, unprotected variable access), and consistency errors (spawn/join mismatch, unreachable statements) before the more expensive verification step.


% ── 5. 为什么还需要形式化验证 ────────────────────────────
% 如果逻辑逻辑模型（LLM）对并发的理解足以生成循环死锁（Cir），为什么不直接要求它们验证正确性呢？答案在于并发性错误的本质。跨线程错误源于独立正确操作之间的组合交互。考虑三个线程以不同的顺序获取三个锁：每个线程的代码在局部都是正确的，但系统却出现了循环死锁。交错操作的数量随着并发实体和同步操作的数量呈指数级增长。LLM本质上是顺序推理器。它们可以理解局部模式（“这个锁在那个锁之前获取”）和常见惯用法（“条件变量 wait 应该放在 while 循环中”），但它们无法系统地枚举 10^6 种交错操作来找到导致死锁的那一种。这正是形式化验证方法的作用：它们提供给定抽象范围内状态空间的“穷举覆盖”。
\subsection{Why LLMs Alone Are Not Enough}
\label{subsec:why-verify}

If LLMs understand concurrency well enough to generate \textsc{Cir}, why not simply ask them to also verify correctness? The answer lies in the nature of concurrency bugs. Cross-thread bugs emerge from \emph{combinatorial interactions} among independently correct operations. Consider three threads acquiring three locks in different orders: each thread's code is locally correct, but the system exhibits circular deadlock. The number of interleavings grows exponentially with the number of concurrent entities and synchronization operations. LLMs are fundamentally sequential reasoners. They can understand local patterns (``this lock is acquired before that one'') and common idioms (``condvar wait should be in a while loop''), but they cannot systematically enumerate $10^6$ interleavings to find the one that leads to deadlock. This is precisely what formal verification methods do: they provide \emph{exhaustive coverage} of the state space within a given abstraction.

This leads to a natural separation of concerns: the LLM is tasked with the semantic heavy-lifting—interpreting the program's intent to generate a precise, alias-free \textsc{Cir}; the formal engine then takes over to rigorously explore all possible execution interleavings for bugs; and finally, the LLM re-enters the loop to diagnose and rectify the specification based on the formal engine’s findings.

% ── 6. 为什么选 Petri 网，为什么是 CVN ───────────────────
% 形式化验证后端的选择并非显而易见。原则上，一些成熟的框架都可以作为分析引擎。在阐述我们选择的理由之前，我们会对主要的替代方案进行比较。显式状态模型检测（例如 SPIN）要求将程序重新表达为领域特定的建模语言（Promela）。虽然功能强大，但这引入了一个手动转换步骤，而我们的架构旨在实现这一步骤的自动化。此外，生成的模型也缺乏用于模块化分析的自然组合结构。有界模型检测（例如 CBMC）将有界执行编码为 SAT/SMT 问题。它对于顺序程序和有界线程数的程序是有效的，但同步原语（条件变量语义、通道阻塞、信号量计数）的编码需要大量的人工工作，并且无法自然地支持需要循环检测的活性属性（饥饿、活锁）。动态分析（例如 Miri、Loom）探索具体的执行过程。Miri 每次执行运行一个调度； Loom 系统地探索了小型测试框架的有限调度方案。两者都不能提供全面的测试覆盖，而且都需要可执行代码，而不是抽象的中间表示。
\subsection{Why Petri Nets, and Why \textsc{Cvn}}
\label{subsec:why-petri}
The choice of formal verification back-end is not obvious. Several established frameworks could, in principle, serve as the analysis engine. We compare the main alternatives before justifying our choice. Explicit-state model checking (e.g., SPIN) requires programs to be re-expressed in a domain-specific modeling language (Promela). While powerful, this imposes a manual translation step that our architecture seeks to automate. The resulting models also lack a natural compositional structure for modular analysis. Bounded model checking (e.g., CBMC) encodes bounded executions as SAT/SMT problems. It is effective for sequential programs and bounded thread counts, but the encoding of synchronization primitives (condvar semantics, channel blocking, semaphore counting) requires substantial manual effort and does not naturally support liveness properties (starvation, livelock) that require cycle detection. Dynamic analysis (e.g., Miri, Loom) explores concrete executions. Miri runs one schedule per execution; Loom systematically explores bounded schedules for small test harnesses. Neither provides exhaustive coverage, and both require executable code rather than an abstract intermediate representation.

Petri nets offer a natural fit for this task for several compelling reasons. First, they provide a native concurrency model where tokens in places represent thread positions and resource availability, and transitions represent operations; concurrency is inherent rather than encoded atop a sequential framework. Second, they boast well-studied decidability properties—deadlock detection, boundedness, and coverability are decidable for P/T nets, while liveness can be verified through strongly connected component analysis on the reachability graph. Third, their compositional structure allows separate subnets for different functions to compose naturally via shared resource places, mirroring the structure of \textsc{Cir}. Finally, their visual interpretability provides a natural graphical representation that greatly aids in debugging and understanding counterexamples.

However, standard P/T nets lack the ability to model data-dependent branching. While Classical Colored Petri Nets (\textsc{CPNs}) address this by attaching typed data to tokens, they introduce significant complexity regarding color sets, bindings, and arc inscriptions. For our setting, this complexity is unnecessary.

% 由于 CIR 已通过全局命名解析了所有资源标识，因此令牌无需携带标识信息。LLM 已将有界并发实例展开为显式的多个 spawn，为每个实例赋予其自身的控制位置命名空间。因此，黑色令牌（不可区分且未着色）足以跟踪线程位置和资源可用性。为了在不使用着色令牌的情况下对共享变量状态和数据相关的分支进行建模，我们引入了一个具有三值求值的全局变量存储 $V$。这引出了我们的第二个形式化模型：并发验证网 (CVN)。CVN 是一个具有特定增强功能的加权 P/T Petri 网。它包含一个全局变量存储 $V$ : \mathsf{VarName} 到 \mathsf{Val}，其中 \mathsf{Val} 包含具体值和一个吸收未知元素 \topval。它利用输入弧上的守卫，这些输入弧是 $V$ 上的布尔表达式，只有当其值为非假时，转换才能启用。它定义了输出弧上的更新，这些更新是在触发时对 $V$ 执行的赋值操作。关键在于，它采用了三值求值：真（启用）、假（禁用）和未知（通过过度近似启用）。这种设计实现了并发验证所需的建模能力——特别是数据相关的分支和变量保护的同步——而无需使用彩色标记。表~\ref{tab:comparison} 总结了这些模型之间的比较。
Because \textsc{Cir} has already resolved all resource identities through global naming, tokens do not need to carry identity information. The LLM has already unfolded bounded concurrent instances into explicit multiple spawns, giving each instance its own control-place namespace. Black tokens—indistinguishable and uncolored—are therefore sufficient to track thread positions and resource availability. To model shared variable state and data-dependent branching without colored tokens, we introduce a global variable store $V$ with three-valued evaluation. This leads to our second formalism: the Concurrency Verification Net (\textsc{Cvn}). \textsc{Cvn} is a weighted P/T Petri net augmented with specific features. It incorporates a global variable store $V$: $\mathsf{VarName}$ to $\mathsf{Val}$, where $\mathsf{Val}$ includes concrete values and an absorbing unknown element $\top$. It utilizes guards on input arcs, which are Boolean expressions over $V$ that must evaluate to non-false for a transition to enable. It defines updates on output arcs, which are assignments to $V$ executed upon firing. Critically, it employs three-valued evaluation: $\mathsf{true}$, $\mathsf{false}$, and $\mathsf{Unknown}$ enabled via over-approximation. This design achieves the necessary modeling power for concurrency verification—specifically data-dependent branching and variable-guarded synchronization—without the overhead of colored tokens. Table~\ref{tab:comparison} summarizes the comparison between these models.

\begin{table}[t]
\caption{Comparison of verification approaches.}
\label{tab:comparison}
\centering\footnotesize
\begin{tabular}{@{}p{1.4cm}ccccc@{}}
\toprule
& \rotatebox{65}{Exhaustive}
& \rotatebox{65}{No manual model}
& \rotatebox{65}{Liveness}
& \rotatebox{65}{Data guards}
& \rotatebox{65}{Simplicity} \\
\midrule
SPIN        & \checkmark &            & \checkmark & \checkmark &            \\
CBMC        & bounded    &            &            & \checkmark &            \\
Miri/Loom   &            & \checkmark &            & \checkmark & \checkmark \\
CPN         & \checkmark &            & \checkmark & \checkmark &            \\
\textbf{\textsc{Cvn}} & \checkmark & \checkmark & \checkmark & \checkmark & \checkmark \\
\bottomrule
\end{tabular}
\vspace{-6pt}
\end{table}

% ── 7. 反例驱动的修复闭环 ────────────────────────────────
% Cir 生成和 Cvn 验证并非独立进行，它们共同构成一个闭环的生成-验证-修复循环。在 Cvn 状态空间搜索过程中发现的错误会生成一个结构化的反例。该反例明确指出错误类型——死锁、信号丢失、饥饿、活锁或通道阻塞。它提供了一个转换跟踪。该跟踪包含每个步骤的 Cir 语句 ID，将抽象执行与具体的 Cir 语句关联起来。它捕获最终状态，指示哪些线程持有哪些锁以及哪些线程正在等待哪些资源。它还包含一个模板修复建议。该建议源自错误类型，例如，统一死锁的锁获取顺序，或使用 while 循环保护信号丢失的等待操作。
\subsection{Closing the Loop: Counterexample-Guided Repair}
\label{subsec:intro-repair}

\textsc{Cir} generation and \textsc{Cvn} verification do not stand alone. They serve the closed generate–verify–repair loop. A bug found during \textsc{Cvn} state-space search produces a structured counterexample. This counterexample specifies the bug kind—deadlock, signal loss, starvation, livelock, or channel blocking. It provides a transition trace. This trace contains the \textsc{Cir} statement ID of every step, anchoring abstract execution to concrete \textsc{Cir} statements. It captures the final state, indicating which threads hold which locks and which threads wait on which resources. It includes a template repair suggestion. This suggestion is derived from the bug kind, such as unifying lock acquisition order for deadlock or protecting a wait with a while loop for signal loss.

% 该结构化的反例会被格式化为修复提示并发送到 LLM。由于反例锚定于 $\mathsf{sid}$s 而非模糊的自然语言描述，LLM 可以生成有针对性的修复——修改特定的 \textsc{Cir} 语句，而不是重写整个程序。修复后的 \textsc{Cir} 通过同一流程重新验证，该过程最多重复 K 轮。这种架构利用了互补性：LLM 提供语义理解（程序应该做什么），而形式引擎提供穷举覆盖（程序在所有交错情况下实际做什么）。两者单独都不够：LLM 无法枚举交错情况，形式引擎也无法理解业务意图。它们共同构成了一个各自都无法独立完成的闭环。
This structured counterexample is formatted into a repair prompt and sent to the LLM. Because the counterexample is anchored to $\mathsf{sid}$s rather than vague natural-language descriptions, the LLM can generate a \emph{targeted} fix---modifying specific \textsc{Cir} statements rather than rewriting the entire program. The repaired \textsc{Cir} is re-validated through the same pipeline, and the process repeats for at most $K$ rounds. This architecture exploits a complementarity: the LLM contributes \emph{semantic understanding} (what the program should do), while the formal engine contributes \emph{exhaustive coverage} (what the program actually does under all interleavings). Neither alone is sufficient: the LLM cannot enumerate interleavings, and the formal engine cannot understand business intent. Together, they close a loop that neither can close independently.

% ── 8. 贡献 ──────────────────────────────────────────────
\subsection{Contributions}
\label{subsec:contributions}

This paper makes the following contributions:

\begin{enumerate}
\item We show that the alias barrier can be bypassed entirely by exploiting LLM semantic understanding to generate an alias-free intermediate representation. The hard question ''what aliases?'' is replaced by an easy question ``what should be shared?'' that LLMs answer reliably.

\item \textsc{Cir} features global resource naming, explicit protection mappings, and statement-level anchors. A static checker enforces 61~well-formedness rules in 9~categories.

\item \textsc{Cvn} extends weighted P/T nets with global-variable guards and three-valued evaluation, achieving the modeling power of colored Petri nets for concurrency analysis without colored-token complexity.

\item A faithful translation from \textsc{Cir} to \textsc{Cvn} with a step-correspondence theorem guaranteeing that every \textsc{Cir} execution has a corresponding \textsc{Cvn} firing sequence and vice versa.

\item A two-layer verification architecture combining fast \textsc{Cir} static checking with exhaustive \textsc{Cvn} state-space search, producing structured counterexamples for LLM-driven repair.

\item A test matrix of 10~representative patterns demonstrating that the pipeline catches cross-thread bugs that neither LLMs alone nor existing static analyzers can reliably detect.
\end{enumerate}

Section~\ref{sec:motivation} presents a motivating example.
Section~\ref{sec:cir} defines \textsc{Cir}.
Section~\ref{sec:cvn} defines \textsc{Cvn}.
Section~\ref{sec:translation} presents the translation.
Section~\ref{sec:analysis} describes state-space analysis.
Section~\ref{sec:repair} covers counterexamples and repair.
Section~\ref{sec:loops} discusses loop modeling.
Section~\ref{sec:formal} states formal properties.
Section~\ref{sec:tests} presents the test matrix.
Section~\ref{sec:related} surveys related work.
Section~\ref{sec:conclusion} concludes.


\section{Motivation}
\label{sec:motivation}

\subsection{Running Example: A Lost Wake-up Bug}
\label{subsec:running-example}

Figure~\ref{fig:running-rust} presents a typical synchronization defect in Rust involving a condition variable. In the \texttt{notifier} function, the notification (\texttt{notify\_one}) is invoked before the readiness flag is updated. If the \texttt{notifier} completes this sequence before the \texttt{worker} thread enters the \texttt{wait} state, the signal is lost, causing the \texttt{worker} to block indefinitely. Despite this critical liveness failure, the program adheres to Rust's ownership rules and compiles without warnings.

\begin{figure}[h]
\begin{lstlisting}[language=Rust, basicstyle=\small\ttfamily, frame=single]
fn worker(shared: Arc<(Mutex<bool>, Condvar)>) {
let (m, cv) = &*shared;
let mut g = m.lock().unwrap();
while !*g {
g = cv.wait(g).unwrap();
}
}

fn notifier(shared: Arc<(Mutex<bool>, Condvar)>) {
let (m, cv) = &*shared;
let mut g = m.lock().unwrap();
cv.notify_one();   // BUG: notify before flag update
*g = true;         // too late
}
\end{lstlisting}
\caption{A condition-variable signal-loss bug in Rust.}
\label{fig:running-rust}
\end{figure}

\subsection{The Complexity of Conventional Static Analysis}
\label{subsec:static-barriers}

Detecting this bug through traditional static analysis requires the simultaneous resolution of three interdependent challenges: identity of shared resources, implicit protection invariants, and non-deterministic interleaving. First, establishing the \textbf{identity of shared resources} requires proving that separate threads operate on the exact same synchronization primitives. In Rust, this necessitates tracing values through heap allocation, \texttt{Arc} cloning, and closure capture. Such pointer-aliasing problems frequently lead to a compromise between precision and scalability, often resulting in either \emph{false positives} due to over-approximation or \emph{false negatives} from under-approximation. 

Beyond the aliasing hurdle, implicit protection invariants such as the logical guarding of a boolean flag by a specific mutex are typically buried within implementation patterns rather than being explicitly declared. Traditional tools are often tightly coupled to specific language semantics, such as Rust’s \texttt{MutexGuard}, which restricts their generalizability across different memory models or languages. Even if these resource identities and invariants are successfully inferred, the analyzer must still perform an exhaustive interleaving analysis to expose the specific schedule where a notification is lost. Existing static analyzers frequently struggle at the initial stages of this pipeline. By shifting the focus from language-specific implementation artifacts to a language-agnostic abstraction of concurrent logic, we decouple the verification of synchronization patterns from the inherent complexities of pointer aliasing and source-level syntax.

\subsection{Bypassing Aliasing via \textsc{Cir} Abstraction}
\label{subsec:cir-abstraction}

The \textsc{Cir} in Figure~\ref{fig:running-cir} reframes the problem by shifting from implementation level inference to design level specification. Our architecture leverages an LLM to synthesize this alias-free representation from the program intent.

\begin{figure}[h]
\begin{lstlisting}[basicstyle=\small\ttfamily, frame=single]
resources:
m0:    { kind: Mutex, mode: Sync }
cv0:   { kind: Condvar, paired_with: m0 }
ready: { kind: Var, type: Bool, init: false }
protection:
ready: [m0]
functions:
worker:
body:
- { sid: w1, op: {lock: m0}, transfer: {next: w2} }
- { sid: w2, op: {read: ready}, transfer: {branch: {...}} }
- { sid: w3, op: {wait: {cv: cv0, mutex: m0}}, ... }
notifier:
body:
- { sid: n1, op: {lock: m0}, ... }
- { sid: n2, op: {notify_one: cv0}, transfer: {next: n3} }
- { sid: n3, op: {write: {var: ready, val: true}}, ... }
\end{lstlisting}
\caption{\textsc{Cir} for the running example, featuring global naming and explicit protection.}
\label{fig:running-cir}
\end{figure}

In \textsc{Cir}, shared identities are established by construction through global naming (\texttt{m0}, \texttt{cv0}). The protection relationship is elevated to a first-class declaration. By stripping away pointer indirection and heap-tracking, \textsc{Cir} transforms an intractable alias-analysis problem into a direct verification task on a clean, concurrent backbone.

\subsection{Synergistic Discovery and Repair}
\label{subsec:closing-loop}

The \textsc{Cir} is translated into a \textsc{Cvn}, which explores all interleavings. The formal engine identifies that if the transition at \texttt{n2} fires while the worker has not reached the wait place at \texttt{w3}, a signal loss occurs.

Crucially, the resulting counterexample is structured and anchored to specific statement IDs (\texttt{sid}). This allows the LLM to perform a targeted repair (e.g., swapping \texttt{n2} and \texttt{n3}) rather than a blind rewrite. This synergy exploits a fundamental division of labor: the LLM provides the semantic intuition to resolve resource identities, while the formal engine provides the exhaustive coverage required to expose edge-case interleavings.




% !TeX root = ../main.tex
\section{System Architecture}
\label{sec:architecture}

% 该架构由LLM驱动的Cir生成阶段、确定性验证层、机械翻译阶段和本地验证修复循环组成.由此产生的流水线追踪从自然语言规范到可接受的Cir工件的信息流,同时将模型质量检查与并发行为分析分离.
The architecture consists of an LLM-driven \textsc{Cir} generation stage, a deterministic validation layer, a mechanical translation stage, and a local verification-and-repair loop. The resulting pipeline traces information flow from natural-language specification to accepted \textsc{Cir} artifact while separating model-quality checking from concurrency-behavior analysis.

\subsection{Pipeline Overview}
\label{subsec:pipeline-overview}

% 图展示完整闭环,从规格到 CIR,到 CVN,到局部修复.
Figure~\ref{fig:pipeline} summarizes the end-to-end pipeline.

\begin{figure}
\centering
\includegraphics[width=0.45\columnwidth]{figures/graphviz.png}
\caption{The generate-verify-repair pipeline.}
\Description{A five-stage pipeline. A natural-language specification is turned into a CIR artifact by an LLM. The CIR passes a static checker, is translated into a CVN model, and is explored by state-space analysis. If a bug is found, structured feedback is sent to the LLM for repair; otherwise goal reachability is checked before acceptance.}
\label{fig:pipeline}
\end{figure}

% Step1: LLM 读取自然语言需求或系统规范,并生成 Cir 工件.当前原型使用 Rust 示例进行评估,但此阶段的输出语言是 Cir,而非 Rust 语法.LLM 的任务是恢复预期的共享资源、保护映射和语句级控制流,而不是生成可执行程序.
\emph{Stage 1: \textsc{Cir} generation.}
An LLM reads a natural-language requirement or system specification and produces a \textsc{Cir} artifact. The current prototype is evaluated on Rust examples, but the output language of this stage is \textsc{Cir}, not Rust syntax. The LLM is asked to recover the intended shared resources, protection mappings, and statement-level control flow rather than to emit an executable program.

% Step2: 确定性检查器会根据分为 9 个类别的 61 条格式良好性规则来验证生成的 \textsc{Cir}.此阶段可以及早发现格式错误的工件:缺少字段、未定义的资源、不兼容的操作、损坏的 spawn/join 结构以及保护映射错误.
\emph{Stage 2: \textsc{Cir} static checking.}
A deterministic checker validates the generated \textsc{Cir} against 61 well-formedness rules organized in 9 categories. This stage catches malformed artifacts early: missing fields, undefined resources, incompatible operations, broken spawn/join structure, and protection-mapping errors.

% Step3: 机械翻译器将验证后的 Cir 转换为 Cvn 模型.翻译是直接的:每个 Cir 语句产生至少一个锚定变迁,不进行优化或变迁合并.这使得后续的正确性论证简单,并保留诊断可追踪性.
\emph{Stage 3: Translation to \textsc{Cvn}.}
A stateless translator converts the validated \textsc{Cir} into a \textsc{Cvn}. The translation is intentionally direct: each \textsc{Cir} statement gives rise to at least one anchored transition, and no optimization or transition merging is performed. This keeps the later correctness argument simple and preserves diagnostic traceability.

% Step4: 本地状态空间分析通过穷举可达性和基于 SCC 的检查来执行.此阶段不消耗 LLM 提示预算.其角色是找到在显式足够的验证模型中存在的 bug 见证,该模型足够紧凑,可以保持有界.
\emph{Stage 4: Local state-space analysis.}
The \textsc{Cvn} is analyzed locally through exhaustive reachability and SCC-based checks \cite{Holzmann1997SPIN,Flanagan2005DPOR,Musuvathi2008FairStateless,Esparza2008Unfoldings}. This stage does not consume LLM prompt budget. Its role is to find bug witnesses in a verification model that is explicit enough for formal reasoning but compact enough to remain bounded.

% Step5: 结构化修复反馈.如果分析发现 bug,验证器不会返回原始网或原始标记转储.相反,它打包一个与 Cir 语句标识符锚定的修复导向诊断,并将其紧凑的报告发送回 LLM.修正后的 Cir 然后重新进入管道第 2 阶段.如果没有发现 bug,则使用相同的可达状态集来检查业务目标可达性,然后进行接受.
\emph{Stage 5: Structured repair feedback.}
If analysis finds a bug, the verifier does not return a raw net or a raw marking dump. Instead, it packages a repair-oriented diagnostic anchored to \textsc{Cir} statement identifiers and sends that compact report back to the LLM. The corrected \textsc{Cir} then re-enters the pipeline at Stage~2. If no bug is found, the same reachable state set is reused to check business-goal reachability before acceptance.

\subsection{Verification and Repair Hierarchy}
\label{subsec:two-layer}
\label{subsec:repair-tier}

% 两层验证体系与三层修复响应机制相结合.验证层将模型质量与并发行为区分开来.修复层则决定在发现问题后需要采取何种程度的干预措施.
A two-layer verification hierarchy is coupled with three repair responses. The verification layers separate model quality from concurrency behavior. The repair tiers determine how much intervention is needed once a problem is found.

\emph{Layer 1: \textsc{Cir} well-formedness checking.}
This layer targets generation artifacts. Since \textsc{Cir} is produced by an LLM, it may contain shallow structural or semantic inconsistencies even when the intended design is clear. Layer~1 serves as a deterministic filter that checks whether the artifact is internally coherent and ready for formal analysis.

\emph{Layer 2: \textsc{Cvn} state-space analysis.}
This layer assumes that the \textsc{Cir} is already well formed and considers whether the modeled concurrency logic admits a bad reachable behavior. The bugs found here are interleaving-dependent and therefore require exhaustive analysis rather than local pattern recognition.

\emph{Tier 1: deterministic auto-fix.}
If Layer~1 reports a trivial, local, and unambiguous issue, the checker repairs it directly using pre-defined rewrites. This avoids spending LLM rounds on mechanical corrections.

\emph{Tier 2: semantic re-generation.}
If Layer~1 reports a more substantial structural problem, the verifier sends the diagnostic back to the LLM for a localized regeneration of the affected \textsc{Cir} fragment.

\emph{Tier 3: logic-driven repair.}
If Layer~2 finds a concurrency violation, the verifier constructs a structured diagnostic from the bug witness and asks the LLM to revise the affected synchronization logic. This is the most expensive repair tier and the one that carries the main methodological novelty of the paper.

% 表格保留,但把内容收束成"谁发现什么、怎么修".
Table~\ref{tab:two-layer} summarizes this division of responsibilities.

\begin{table}[t]
\caption{Verification and repair hierarchy.}
\label{tab:two-layer}
\centering\footnotesize
\begin{tabular}{@{}lll@{}}
\toprule
\textbf{Issue} & \textbf{Detected by} & \textbf{Primary response} \\
\midrule
Missing \texttt{drop}         & Layer 1 (\textsc{Cir}) & Tier 1 auto-fix \\
Undefined resource            & Layer 1 (\textsc{Cir}) & Tier 2 LLM repair \\
Spawn/join mismatch           & Layer 1 (\textsc{Cir}) & Tier 2 LLM repair \\
Broken protection mapping     & Layer 1 (\textsc{Cir}) & Tier 2 LLM repair \\
Two-lock deadlock             & Layer 2 (\textsc{Cvn}) & Tier 3 LLM repair \\
Condvar signal loss           & Layer 2 (\textsc{Cvn}) & Tier 3 LLM repair \\
Three-lock circular deadlock  & Layer 2 (\textsc{Cvn}) & Tier 3 LLM repair \\
Writer starvation             & Layer 2 (\textsc{Cvn}) & Tier 3 LLM repair \\
Channel-mutex cross-deadlock  & Layer 2 (\textsc{Cvn}) & Tier 3 LLM repair \\
\bottomrule
\end{tabular}
\end{table}

% 这种层级结构有两个目的.首先,它可以避免将耗费大量资源的二层分析浪费在那些本应更早被拒绝的畸形工件上.其次,它可以确保层级模型管理在合适的抽象层级上接收反馈:当模型畸形时提供结构性反馈,当并发逻辑错误时提供基于执行的反馈,只有在确认模型无缺陷后才提供目标违背反馈.
This hierarchy serves two purposes. First, it prevents expensive Layer~2 analysis from being wasted on malformed artifacts that should have been rejected much earlier. Second, it ensures that the LLM receives feedback at the right level of abstraction: structural feedback when the model is malformed, execution-grounded feedback when the concurrency logic is wrong, and goal-violation feedback only after bug-freedom has been established.


\section{Formal Definitions}
\label{sec:formal-defs}

This section defines the three core formalisms: the Concurrency Intermediate Representation (\textsc{Cir}), the Concurrency Verification Net (\textsc{Cvn}), and the translation from
\textsc{Cir} to \textsc{Cvn}. Each subsection first states the formal structure, then describes its components, and finally presents the associated error taxonomy.

% ══════════════════════════════════════════════════════════════
%  4.1  CIR
% ══════════════════════════════════════════════════════════════
\subsection{CIR: Concurrency Intermediate Representation}
\label{subsec:cir-def}

\subsubsection{Top-Level Structure}

\begin{definition}[\textsc{Cir} artifact]
\label{def:cir}
A \textsc{Cir} artifact is a five-tuple
\[
  \mathcal{I} = ( R,\; G,\; F,\; \Sigma,\; e )
\]
where
\begin{itemize}[leftmargin=1.4em]
  \item $R$ is a finite resource table,
  \item $G$ is a protection mapping,
  \item $F$ is a finite set of function definitions,
  \item $\Sigma$ is a partial function-summary map, and
  \item $e \in F$ is the entry function.
\end{itemize}
\end{definition}

\subsubsection{Resource Table}

The resource table $R$ maps each resource name to a kind and a configuration. Resources are classified into two families.

\textbf{Synchronization primitives.} Each primitive has a kind and a mode.
\[
\begin{aligned}
\mathit{kind} &\in \{ \textsf{Mutex},\; \textsf{RwLock},\;
  \textsf{Condvar},\; \textsf{Semaphore}(n),\; \textsf{Channel} \} \\
\mathit{mode} &\in \{ \textsf{Sync},\; \textsf{Async} \}
\end{aligned}
\]
A \textsf{Condvar} carries a \texttt{paired\_with} field naming the associated mutex. A \textsf{Semaphore}$(n)$ carries an initial permit count $n \geq 1$.

\textbf{Shared variables.} Each variable has a kind, a base type, and an initial value.
\[
\begin{aligned}
\mathit{kind} &\in \{ \textsf{Var},\; \textsf{Atomic} \} \\
\mathit{type} &\in \{ \textsf{Bool},\; \textsf{Int},\;
  \textsf{Float},\; \textsf{String},\; \textsf{Enum}(S) \} \\
\mathit{init} &\in \mathit{Values}(\mathit{type})
\end{aligned}
\]
where $\textsf{Enum}(S)$ denotes an enumeration over a finite set of variants $S$.

Every resource has a unique global name. There are no pointers, no heap indirection, and no aliasing. Functions reference resources directly by name.

\subsubsection{Protection Mapping}

The protection mapping is a partial function
\[
  G : \mathit{VarName} \rightharpoonup \mathcal{P}(\mathit{LockName})
\]
that maps each \textsf{Var} resource to a nonempty set of lock resources. A declaration $G(x) = \{m_1, m_2\}$ means that any access to variable $x$ requires holding at least one of $m_1$ or $m_2$.

\textsf{Atomic} resources do not participate in the protection mapping. They are accessed through atomic operations that provide their own synchronization guarantees.

The protection mapping is used only for \textsc{Cir} static checking. It does not produce any \textsc{Cvn} structure.

\subsubsection{Functions and Statements}

A function definition is a triple $(f, k, B)$ where $f$ is the function name, $k \in \{ \textsf{normal}, \textsf{async}, \textsf{closure} \}$ is the function kind, and $B$ is the function body.

A function body $B$ is a finite set of statements. Each statement is a triple
\[
  s = (\mathit{sid},\; \mathit{op},\; \mathit{transfer})
\]
where $\mathit{sid}$ is a unique identifier within the function, $\mathit{op}$ is the operation, and $\mathit{transfer}$ specifies control flow after the operation.

Functions do not take parameters and do not return values. All inter-function communication occurs through globally named resources.

\subsubsection{Operations}

Table~\ref{tab:cir-ops} lists the complete operation set. Operations are grouped into four categories.

\begin{table}[t]
\caption{\textsc{Cir} operations.}
\label{tab:cir-ops}
\centering\footnotesize
\begin{tabular}{@{}llp{3.2cm}@{}}
\toprule
\textbf{Category} & \textbf{Operation} & \textbf{Target} \\
\midrule
\multirow{2}{*}{Lock}
  & \textsf{lock}, \textsf{drop}
  & Mutex, RwLock \\
  & \textsf{read\_lock}, \textsf{write\_lock}
  & RwLock \\[3pt]
\multirow{3}{*}{Synchronization}
  & \textsf{wait}, \textsf{notify\_one}, \textsf{notify\_all}
  & Condvar \\
  & \textsf{acquire}, \textsf{release}
  & Semaphore \\
  & \textsf{send}, \textsf{recv}
  & Channel \\[3pt]
\multirow{3}{*}{Data}
  & \textsf{read}, \textsf{write}
  & Var \\
  & \textsf{load}, \textsf{store}
  & Atomic \\
  & \textsf{cas}(\mathit{expected}, \mathit{new})
  & Atomic \\[3pt]
\multirow{3}{*}{Control}
  & \textsf{spawn}, \textsf{join}
  & function (OS thread) \\
  & \textsf{spawn\_async}, \textsf{await}
  & function (async task) \\
  & \textsf{call}
  & function \\
  & \textsf{return}
  & --- \\
\bottomrule
\end{tabular}
\end{table}

\textsf{spawn} pairs with \textsf{join} for OS threads. \textsf{spawn\_async} pairs with \textsf{await} for asynchronous tasks. Cross-pairing is a static error. \textsf{acquire} and
\textsf{release} are used for semaphores and are distinct from \textsf{lock} and \textsf{drop}.

\subsubsection{Control Transfer}

The $\mathit{transfer}$ field takes one of the following forms.

\[
\begin{aligned}
\mathit{transfer} ::= \;&\textsf{next}(\mathit{sid'}) \\
  \mid\;& \textsf{branch}(\beta,\; \mathit{sid}_t,\; \mathit{sid}_f) \\
  \mid\;& \textsf{switch}(x,\; [(v_1, \mathit{sid}_1), \ldots, (v_k, \mathit{sid}_k)]) \\
  \mid\;& \textsf{done}
\end{aligned}
\]

\textsf{next} transfers control to a single successor. \textsf{branch} evaluates a boolean condition $\beta$ over variable state and transfers to $\mathit{sid}_t$ if true or $\mathit{sid}_f$ if false. \textsf{switch} dispatches on the value of a variable. \textsf{done} terminates the function.

Control flow is determined by explicit target identifiers in the $\mathit{transfer}$ field. The ordering of statements within the body has no semantic significance. A \textsc{Cir} function body is a labeled directed graph, not a list.

\subsubsection{Branch Conditions}

A branch condition is a boolean expression over global variable state:
\[
\begin{aligned}
\beta ::= \;&(x \;\mathit{cmp}\; v) \\
  \mid\;& \beta_1 \wedge \beta_2 \\
  \mid\;& \neg \beta
\end{aligned}
\]
where $x$ is a variable name, $v$ is a literal value, and $\mathit{cmp} \in \{ =, \neq, <, \leq, >, \geq \}$.

\subsubsection{Function Summary}

The summary map $\Sigma$ provides summaries only for functions whose bodies are not modeled in $F$. A summary is a tuple
\[
  \Sigma(f) = (\mathit{reads},\; \mathit{writes},\;
    \mathit{calls},\; \mathit{has\_concurrency})
\]
where $\mathit{reads}$ and $\mathit{writes}$ are sets of resource names, $\mathit{calls}$ is a set of callee names, and $\mathit{has\_concurrency}$ is a boolean. Fully modeled functions do
not have summaries.

\subsubsection{CIR Error Taxonomy}
\label{subsubsec:cir-errors}

The static checker enforces 61 well-formedness rules organized in 9 categories with prefix E. Table~\ref{tab:cir-errors} summarizes the categories.

\begin{table}[t]
\caption{\textsc{Cir} error categories (61 rules total).}
\label{tab:cir-errors}
\centering\footnotesize
\begin{tabular}{@{}llp{3.6cm}@{}}
\toprule
\textbf{Code} & \textbf{Category} & \textbf{Examples} \\
\midrule
E0xx & Structural
  & Missing fields, invalid format \\
E1xx & Name resolution
  & Undefined resource, duplicate name \\
E2xx & Type
  & Non-boolean branch condition \\
E3xx & Resource compatibility
  & \textsf{lock} on non-lock resource, unprotected variable access \\
E4xx & Concurrency pairing
  & \textsf{spawn} without \textsf{join}, cross-pairing \\
E5xx & Lock safety
  & Missing \textsf{drop}, double lock, inconsistent lock order \\
E6xx & Control flow
  & Unreachable statement, missing \textsf{return} \\
E7xx & Protection mapping
  & \textsf{Atomic} in protection, \textsf{Var} without protection \\
E8xx & Function summary
  & Summary conflicts with body \\
\bottomrule
\end{tabular}
\end{table}

% ══════════════════════════════════════════════════════════════
%  4.2  CVN
% ══════════════════════════════════════════════════════════════
\subsection{CVN: Concurrency Verification Net}
\label{subsec:cvn-def}

\subsubsection{Overview}

A \textsc{Cvn} is a weighted Place/Transition Petri net augmented with a global variable store and three-valued guard evaluation. It is not a colored Petri net. All tokens are indistinguishable black tokens. There are no color sets, no binding expressions, and no
token-carried data.

\subsubsection{Net Structure}

\begin{definition}[\textsc{Cvn}]
\label{def:cvn}
A Concurrency Verification Net is an eight-tuple
\[
  \mathcal{N} = ( P,\; T,\; A_{\mathrm{in}},\; A_{\mathrm{out}},\;
    V,\; I_m,\; I_v,\; \mu )
\]
where
\begin{itemize}[leftmargin=1.4em]
  \item $P = P_c \uplus P_r \uplus P_w$ is a finite set of places, partitioned into control places $P_c$, resource places $P_r$, and wait places $P_w$,
  \item $T$ is a finite set of transitions,
  \item $A_{\mathrm{in}} \subseteq P \times T$ is the set of input arcs, each annotated with a weight $w \in \mathbb{N}^+$ and a guard $g$,
  \item $A_{\mathrm{out}} \subseteq T \times P$ is the set of output arcs, each annotated with a weight $w \in \mathbb{N}^+$ and an update $u$,
  \item $V : \mathit{VarName} \to \mathit{Val}$ is the global variable store,
  \item $I_m : P \to \mathbb{N}_0$ is the initial marking,
  \item $I_v$ is the initial variable valuation, and
  \item $\mu : T \rightharpoonup \mathit{SID}^+$ is the anchor mapping from transitions to \textsc{Cir} statement identifiers (optional, controlled by a feature flag).
\end{itemize}
\end{definition}

\subsubsection{Place Families}

\textbf{Control places} $P_c$. There is one control place $\mathit{cp}(f, \mathit{sid})$ for each statement in each function. A token in $\mathit{cp}(f, \mathit{sid})$ means that a thread is
at statement $\mathit{sid}$ in function $f$. Each function also has a return place $\mathit{cp}(f, \mathit{ret})$ that receives a token when the function returns.

\textbf{Resource places} $P_r$. There is one resource place $\mathit{rp}(r)$ for each synchronization primitive. The initial token count depends on the resource kind (Table~\ref{tab:init-tokens}). \textsf{Var} and \textsf{Atomic} resources do not have resource
places. Their state is stored in the variable store $V$.  Each condition variable $\mathit{cv}$ with associated wait sites $\{w_1, \ldots, w_k\}$ introduces the following \textsc{Cvn}
elements.

\textbf{Places.}
\begin{itemize}[leftmargin=1.4em]
  \item $\mathit{rp}(\mathit{cv})$: a resource place for the condition variable itself, initialized with 0 tokens. A \textsf{notify\_one} deposits one token; a \textsf{wait} wake-up consumes one token.
  \item $\mathit{wp}(w_i)$: a wait place for each wait site $w_i$. A token here means a thread is blocked at $w_i$.
  \item $\mathit{ra}(w_i)$: a reacquire place for each wait site $w_i$. A token here means a thread has been woken and is waiting to reacquire the associated mutex.
\end{itemize}

\textbf{Global variables.}
\begin{itemize}[leftmargin=1.4em]
  \item $\mathit{nw}_{\mathit{cv}}$: an integer variable counting
        the number of threads currently waiting on $\mathit{cv}$.
        Initialized to 0. Incremented when a thread enters a wait
        place, decremented when a thread leaves a wait place.
  \item $\mathit{na}_{w_i}$: a boolean variable for each wait site
        $w_i$. Initialized to $\mathsf{false}$. Set to
        $\mathsf{true}$ by \textsf{notify\_all}, reset to
        $\mathsf{false}$ when the thread enters the wait place or
        wakes up via this flag.
\end{itemize}

The waiter count $\mathit{nw}_{\mathit{cv}}$ enables the
\textsf{notify} transitions to determine whether any thread is
waiting, using a guard on a global variable rather than inspecting
token counts in places. The per-site flag $\mathit{na}_{w_i}$
provides a second wake-up channel for \textsf{notify\_all} that does
not consume a token from $\mathit{rp}(\mathit{cv})$.

\begin{table}[t]
\caption{Initial token counts for resource places.}
\label{tab:init-tokens}
\centering\footnotesize
\begin{tabular}{@{}ll@{}}
\toprule
\textbf{Resource kind} & \textbf{Initial tokens in $\mathit{rp}(r)$} \\
\midrule
\textsf{Mutex}          & 1 \\
\textsf{RwLock}         & $N$ (total concurrent entities) \\
\textsf{Semaphore}$(n)$ & $n$ \\
\textsf{Channel}        & 0 \\
\bottomrule
\end{tabular}
\end{table}

\textbf{Wait places} $P_w$. There is one wait place $\mathit{wp}(\mathit{sid})$ for each \textsf{wait} statement in the \textsc{Cir}. A token in $\mathit{wp}(\mathit{sid})$ means that
a thread is blocked at that wait site.

\subsubsection{Value Domain}

\begin{definition}[Value domain]
\label{def:val}
The value domain is
\[
  \mathit{Val} = \mathit{Concrete}(\tau) \;\cup\; \{ \top \}
\]
where $\tau$ ranges over the base types (\textsf{Bool}, \textsf{Int}, \textsf{Float}, \textsf{String}, \textsf{Enum}). The element $\top$ (unknown) is an absorbing element: once a variable takes the value $\top$, no operation can restore it to a concrete value.
\end{definition}

Expressions over the value domain come in two forms.

Value expressions:
\[
  \mathit{Expr} ::= \mathit{Lit}(v) \;\mid\;
    \mathit{Ref}(x) \;\mid\;
    \mathit{BinOp}(\mathit{Expr}, \oplus, \mathit{Expr})
\]

Boolean expressions:
\[
  \mathit{BoolExpr} ::= \textsf{True} \;\mid\;
    \mathit{Cmp}(\mathit{Expr}, \mathit{cmp}, \mathit{Expr}) \;\mid\;
    \mathit{And}(\mathit{BoolExpr}, \mathit{BoolExpr}) \;\mid\;
    \mathit{Not}(\mathit{BoolExpr})
\]

\subsubsection{Three-Valued Guard Evaluation}

\begin{definition}[Guard evaluation]
\label{def:guard-eval}
A guard $g$ is a $\mathit{BoolExpr}$. Given a variable store $V$, the evaluation $\mathit{eval}(g, V)$ returns one of three results:
\begin{itemize}[leftmargin=1.4em]
  \item $\mathsf{true}$: the guard is satisfied.
  \item $\mathsf{false}$: the guard is not satisfied.
  \item $\mathsf{unknown}$: the guard involves a variable whose
        value is $\top$.
\end{itemize}
When any operand of a comparison is $\top$, the comparison evaluates to $\mathsf{unknown}$. The conjunction $\mathsf{unknown} \wedge \mathsf{false}$ evaluates to $\mathsf{false}$. The conjunction $\mathsf{unknown} \wedge \mathsf{true}$ evaluates to $\mathsf{unknown}$.
\end{definition}

The evaluation semantics for $\mathsf{unknown}$ is an over-approximation. When a guard evaluates to $\mathsf{unknown}$, the transition is treated as enabled. If the guard arises from a branch statement, both the true-branch and the false-branch transitions are enabled simultaneously. The state-space search explores both paths. This ensures that no reachable state is missed
at the cost of potentially exploring infeasible states.

Guards are semantically meaningful only on transitions generated from \textsf{branch} and \textsf{switch} statements. All other transitions have guard $\textsf{True}$.

\subsubsection{Enabling and Firing}

\begin{definition}[Enabling]
\label{def:enabling}
A transition $t \in T$ is enabled at state $(M, V)$ if and only if:
\begin{enumerate}[leftmargin=1.4em]
  \item for every input arc $(p, t)$ with weight $w$: $M(p) \geq w$, and
  \item for every input arc $(p, t)$ with guard $g$: $\mathit{eval}(g, V) \neq \mathsf{false}$.
\end{enumerate}
\end{definition}

\begin{definition}[Firing]
\label{def:firing}
When an enabled transition $t$ fires at state $(M, V)$, the successor state $(M', V')$ is computed as follows:
\begin{enumerate}[leftmargin=1.4em]
  \item For every input arc $(p, t)$ with weight $w$:
        $M'(p) = M(p) - w$.
  \item For every output arc $(t, p)$ with weight $w$:
        $M'(p) = M(p) + w$ (applied after step 1).
  \item For places not involved in any arc of $t$:
        $M'(p) = M(p)$.
  \item For every output arc $(t, p)$ with update $u$:
        apply $u$ to $V$, yielding $V'$.
  \item Updates on the same transition have pairwise disjoint
        variable domains.
\end{enumerate}
\end{definition}

\subsubsection{State and Finiteness}

\begin{definition}[State]
\label{def:state}
A \textsc{Cvn} state is a pair $(M, V)$ where $M : P \to \mathbb{N}_0$ is the marking and $V : \mathit{VarName} \to \mathit{Val}$ is the variable store. Two states are equal if and only if they agree on both $M$ and $V$.
\end{definition}

\begin{theorem}[Finite state space]
\label{thm:finite}
The reachable state space of a \textsc{Cvn} is finite if every place has a finite token bound and every variable has a finite value domain.
\end{theorem}

\begin{proof}
The marking space is finite because there are finitely many places and each place holds at most a bounded number of tokens. The variable space is finite because each variable takes values in a finite set of concrete values or $\top$, and $\top$ is absorbing. The product of two finite sets is finite.
\end{proof}

Both conditions hold for nets produced by the translation in Section~\ref{subsec:translation}. Token bounds follow from the translation construction: control places are one-safe, resource
places are bounded by the initial token count, and wait places are one-safe. Variable domains are finite because the \textsc{Cir} restricts variable types to finite base types plus $\top$.

\subsubsection{Transition Kinds}

Each transition carries a classification tag $\mathit{kind}$ for debugging and visualization. The tag does not affect the firing semantics. Table~\ref{tab:transition-kinds} lists the classification.

\begin{table}[t]
\caption{\textsc{Cvn} transition kinds.}
\label{tab:transition-kinds}
\centering\footnotesize
\begin{tabular}{@{}ll@{}}
\toprule
\textbf{Category} & \textbf{Kinds} \\
\midrule
Lock        & Lock, Unlock, ReadLock, ReadUnlock \\
Semaphore   & Acquire, Release \\
Channel     & Send, Recv \\
Variable    & VarRead, VarWrite, AtomicLoad, AtomicStore \\
Branch      & BranchTrue, BranchFalse, Switch \\
CAS         & CasSuccess, CasFailure \\
Concurrency & Spawn, Join \\
Call        & Call \\
Condvar     & CondvarWait, CondvarNotify, CondvarNotifyAll, CondvarReacquire \\
\bottomrule
\end{tabular}
\end{table}

\subsubsection{Well-Formedness}

A \textsc{Cvn} is well-formed if it satisfies conditions W1 through W9.

\begin{itemize}[leftmargin=1.4em]
\item[W1.] $P_c$, $P_r$, and $P_w$ are pairwise disjoint.
\item[W2.] Every arc references an existing place and an existing transition.
\item[W3.] Each control place has at most one input arc per transition.
\item[W4.] Updates on a single transition have pairwise disjoint variable domains.
\item[W5.] Branch transitions come in complementary pairs: for each branch, there exist exactly two transitions sharing the same control input place, with guards $g$ and $\neg g$.
\item[W6.] Switch transitions cover all declared cases of the dispatched variable.
\item[W7.] (When the anchor feature is enabled.) Every transition has a nonempty anchor set in $\mu$.
\item[W8.] All arc weights are strictly positive.
\item[W9.] The initial marking is consistent with resource kinds (Table~\ref{tab:init-tokens}).
\end{itemize}

\subsubsection{CVN Error Taxonomy}

Table~\ref{tab:cvn-errors} lists the \textsc{Cvn} error categories
with prefix V.

\begin{table}[t]
\caption{\textsc{Cvn} error categories.}
\label{tab:cvn-errors}
\centering\footnotesize
\begin{tabular}{@{}llp{3.4cm}@{}}
\toprule
\textbf{Code} & \textbf{Category} & \textbf{Examples} \\
\midrule
V0xx & Structural
  & Duplicate ID, dangling arc, zero weight \\
V1xx & Well-formedness
  & Control arc violation, update conflict, missing anchor \\
V2xx & Branch completeness
  & Unpaired branch, incomplete switch \\
V3xx & Analysis
  & Token deficit, state explosion \\
V4xx & Resource semantics
  & Initial tokens inconsistent with kind \\
\bottomrule
\end{tabular}
\end{table}

% ══════════════════════════════════════════════════════════════
%  4.3  Translation
% ══════════════════════════════════════════════════════════════
\subsection{CIR to CVN Translation}
\label{subsec:translation}

\subsubsection{Overview}

The translator is a pure function
\[
  \mathit{translate} : \textsc{Cir} \to \textsc{Cvn}
\]
that takes a validated \textsc{Cir} artifact and produces a \textsc{Cvn}. The translator is stateless and deterministic. It proceeds in three phases.

\textbf{Phase 1: Resource scanning.} For each synchronization primitive in $R$, generate a resource place $\mathit{rp}(r)$ and set its initial token count according to Table~\ref{tab:init-tokens}. Compute the RwLock concurrency parameter $N$ as the total number of concurrent entities (threads and async tasks). For each \textsf{Var} and \textsf{Atomic} resource, initialize the corresponding entry in $V$.

\textbf{Phase 2: Function body translation.} For each function in $F$, for each statement in the body, generate control places, transitions, and arcs according to the rules in Table~\ref{tab:translation-rules}. For each \textsf{wait} statement, also generate a wait place.

\textbf{Phase 3: Summary translation.} For each function referenced by a \textsf{call} statement that has a summary in $\Sigma$ but no body in $F$, generate a single atomic transition. The transition reads no input from resource places (reads have no effect). For each resource in the $\mathit{writes}$ set, the output arc carries an update that sets the variable to $\top$.

\subsubsection{Translation Rules}

Table~\ref{tab:translation-rules} defines the translation of each \textsc{Cir} operation into \textsc{Cvn} structure. The notation uses the following conventions. $\mathit{cp}(f, s)$ denotes the control place for statement $s$ in function $f$. $\mathit{rp}(r)$ denotes the resource place for resource $r$. $\mathit{wp}(s)$ denotes the wait place for wait statement $s$. An arrow $\to$ separates input arcs (consumed tokens and guards) from output arcs (produced tokens and updates). The subscript after the arrow indicates the transition kind.

\begin{table*}[t]
\caption{Translation rules from \textsc{Cir} to \textsc{Cvn} (complete).
Each row shows the \textsc{Cir} pattern and the generated \textsc{Cvn}
structure. Notation: $\mathit{cp}(f,s)$ = control place,
$\mathit{rp}(r)$ = resource place, $\mathit{wp}(s)$ = wait place,
$\mathit{ra}(s)$ = reacquire place.}
\label{tab:translation-rules}
\centering\footnotesize
\begin{tabular}{@{}p{5.6cm}p{11cm}@{}}
\toprule
\textbf{\textsc{Cir} statement} & \textbf{\textsc{Cvn} transitions and arcs} \\
\midrule

% ── Lock / Unlock ────────────────────────────────────
$(\mathit{sid},\; \textsf{lock}(m),\; \textsf{next}(\mathit{sid}'))$
  & One transition $t$ [Lock]:
    input from $\mathit{cp}(f,\mathit{sid})$ weight 1 and
    $\mathit{rp}(m)$ weight 1;
    output to $\mathit{cp}(f,\mathit{sid}')$ weight 1. \\[4pt]

$(\mathit{sid},\; \textsf{drop}(m),\; \textsf{next}(\mathit{sid}'))$
  & One transition $t$ [Unlock]:
    input from $\mathit{cp}(f,\mathit{sid})$ weight 1;
    output to $\mathit{cp}(f,\mathit{sid}')$ weight 1 and
    $\mathit{rp}(m)$ weight 1. \\[4pt]

% ── RwLock ───────────────────────────────────────────
$(\mathit{sid},\; \textsf{read\_lock}(r),\; \textsf{next}(\mathit{sid}'))$
  & One transition $t$ [ReadLock]:
    input from $\mathit{cp}(f,\mathit{sid})$ weight 1 and
    $\mathit{rp}(r)$ weight 1;
    output to $\mathit{cp}(f,\mathit{sid}')$ weight 1. \\[4pt]

$(\mathit{sid},\; \textsf{write\_lock}(r),\; \textsf{next}(\mathit{sid}'))$
  & One transition $t$ [Lock]:
    input from $\mathit{cp}(f,\mathit{sid})$ weight 1 and
    $\mathit{rp}(r)$ weight $N$;
    output to $\mathit{cp}(f,\mathit{sid}')$ weight 1.
    Consumes all $N$ tokens, blocking all readers and writers. \\[4pt]

$(\mathit{sid},\; \textsf{drop}(r),\; \textsf{next}(\mathit{sid}'))$
  \emph{after} \textsf{read\_lock}
  & One transition $t$ [ReadUnlock]:
    input from $\mathit{cp}(f,\mathit{sid})$ weight 1;
    output to $\mathit{cp}(f,\mathit{sid}')$ weight 1 and
    $\mathit{rp}(r)$ weight 1. \\[4pt]

$(\mathit{sid},\; \textsf{drop}(r),\; \textsf{next}(\mathit{sid}'))$
  \emph{after} \textsf{write\_lock}
  & One transition $t$ [Unlock]:
    input from $\mathit{cp}(f,\mathit{sid})$ weight 1;
    output to $\mathit{cp}(f,\mathit{sid}')$ weight 1 and
    $\mathit{rp}(r)$ weight $N$. \\[4pt]

% ── Semaphore ────────────────────────────────────────
$(\mathit{sid},\; \textsf{acquire}(r),\; \textsf{next}(\mathit{sid}'))$
  & One transition $t$ [Acquire]:
    input from $\mathit{cp}(f,\mathit{sid})$ weight 1 and
    $\mathit{rp}(r)$ weight 1;
    output to $\mathit{cp}(f,\mathit{sid}')$ weight 1. \\[4pt]

$(\mathit{sid},\; \textsf{release}(r),\; \textsf{next}(\mathit{sid}'))$
  & One transition $t$ [Release]:
    input from $\mathit{cp}(f,\mathit{sid})$ weight 1;
    output to $\mathit{cp}(f,\mathit{sid}')$ weight 1 and
    $\mathit{rp}(r)$ weight 1. \\[4pt]

% ── Channel ──────────────────────────────────────────
$(\mathit{sid},\; \textsf{send}(c),\; \textsf{next}(\mathit{sid}'))$
  & One transition $t$ [Send]:
    input from $\mathit{cp}(f,\mathit{sid})$ weight 1;
    output to $\mathit{cp}(f,\mathit{sid}')$ weight 1 and
    $\mathit{rp}(c)$ weight 1. \\[4pt]

$(\mathit{sid},\; \textsf{recv}(c),\; \textsf{next}(\mathit{sid}'))$
  & One transition $t$ [Recv]:
    input from $\mathit{cp}(f,\mathit{sid})$ weight 1 and
    $\mathit{rp}(c)$ weight 1;
    output to $\mathit{cp}(f,\mathit{sid}')$ weight 1. \\[4pt]

% ── Variable read (sequential) ───────────────────────
$(\mathit{sid},\; \textsf{read}(x),\; \textsf{next}(\mathit{sid}'))$
  & One transition $t$ [VarRead]:
    input from $\mathit{cp}(f,\mathit{sid})$ weight 1;
    output to $\mathit{cp}(f,\mathit{sid}')$ weight 1.
    No guard, no update. \\[4pt]

% ── Variable read + branch ───────────────────────────
$(\mathit{sid},\; \textsf{read}(x),\; \textsf{branch}(\beta,
  \mathit{sid}_t, \mathit{sid}_f))$
  & Two transitions sharing input from $\mathit{cp}(f,\mathit{sid})$
    weight 1.
    $t_{\top}$ [BranchTrue]: guard $\beta$, output to
    $\mathit{cp}(f,\mathit{sid}_t)$.
    $t_{\bot}$ [BranchFalse]: guard $\neg\beta$, output to
    $\mathit{cp}(f,\mathit{sid}_f)$. \\[4pt]

% ── Variable write ───────────────────────────────────
$(\mathit{sid},\; \textsf{write}(x, v),\; \textsf{next}(\mathit{sid}'))$
  & One transition $t$ [VarWrite]:
    input from $\mathit{cp}(f,\mathit{sid})$ weight 1;
    output to $\mathit{cp}(f,\mathit{sid}')$ weight 1,
    update $V[x] \leftarrow v$. \\[4pt]

% ── Atomic ops ───────────────────────────────────────
$(\mathit{sid},\; \textsf{load}(x),\; \textsf{next}(\mathit{sid}'))$
  & One transition $t$ [AtomicLoad]: same structure as
    \textsf{read}. \\[4pt]

$(\mathit{sid},\; \textsf{store}(x, v),\; \textsf{next}(\mathit{sid}'))$
  & One transition $t$ [AtomicStore]: same structure as
    \textsf{write}, update $V[x] \leftarrow v$. \\[4pt]

$(\mathit{sid},\; \textsf{cas}(x, e, n),\;
  \textsf{next}(\mathit{sid}'))$
  & Two transitions sharing input from $\mathit{cp}(f,\mathit{sid})$.
    $t_s$ [CasSuccess]: guard $V[x] = e$, output to
    $\mathit{cp}(f,\mathit{sid}')$, update $V[x] \leftarrow n$.
    $t_f$ [CasFailure]: guard $V[x] \neq e$, output to
    $\mathit{cp}(f,\mathit{sid}')$, no update. \\[4pt]

% ── Condvar wait ─────────────────────────────────────
% \multicolumn{2}{@{}l}{\textbf{Condvar operations} (cv has
%   paired mutex $m$; wait sites of cv are $\{w_1, \ldots, w_k\}$)} \\[4pt]

$(\mathit{sid},\; \textsf{wait}(\mathit{cv}, m),\;
  \textsf{next}(\mathit{sid}'))$
  & Four transitions.
    \newline
    $t_{\mathit{enter}}$ [CondvarWaitEnter]:
    input from $\mathit{cp}(f,\mathit{sid})$ weight 1;
    output to $\mathit{wp}(\mathit{sid})$ weight 1 and
    $\mathit{rp}(m)$ weight 1;
    update $\mathit{nw}_{\mathit{cv}} \leftarrow
    \mathit{nw}_{\mathit{cv}} + 1$,\;
    $\mathit{na}_{\mathit{sid}} \leftarrow \mathsf{false}$.
    \newline
    $t_{\mathit{wake1}}$ [CondvarWakeByNotify]:
    input from $\mathit{wp}(\mathit{sid})$ weight 1 and
    $\mathit{rp}(\mathit{cv})$ weight 1;
    output to $\mathit{ra}(\mathit{sid})$ weight 1;
    update $\mathit{nw}_{\mathit{cv}} \leftarrow
    \mathit{nw}_{\mathit{cv}} - 1$.
    \newline
    $t_{\mathit{wakeA}}$ [CondvarWakeByNotifyAll]:
    input from $\mathit{wp}(\mathit{sid})$ weight 1;
    guard $\mathit{na}_{\mathit{sid}} = \mathsf{true}$;
    output to $\mathit{ra}(\mathit{sid})$ weight 1;
    update $\mathit{nw}_{\mathit{cv}} \leftarrow
    \mathit{nw}_{\mathit{cv}} - 1$,\;
    $\mathit{na}_{\mathit{sid}} \leftarrow \mathsf{false}$.
    \newline
    $t_{\mathit{reacq}}$ [CondvarReacquire]:
    input from $\mathit{ra}(\mathit{sid})$ weight 1 and
    $\mathit{rp}(m)$ weight 1;
    output to $\mathit{cp}(f,\mathit{sid}')$ weight 1. \\[4pt]

% ── Condvar notify_one ───────────────────────────────
$(\mathit{sid}_n,\; \textsf{notify\_one}(\mathit{cv}),\;
  \textsf{next}(\mathit{sid}_n'))$
  & Two transitions.
    \newline
    $t_{\mathit{notify}}$ [CondvarNotify]:
    input from $\mathit{cp}(f,\mathit{sid}_n)$ weight 1;
    guard $\mathit{nw}_{\mathit{cv}} > 0$;
    output to $\mathit{cp}(f,\mathit{sid}_n')$ weight 1 and
    $\mathit{rp}(\mathit{cv})$ weight 1.
    \newline
    $t_{\mathit{lost}}$ [CondvarNotifyLost]:
    input from $\mathit{cp}(f,\mathit{sid}_n)$ weight 1;
    guard $\mathit{nw}_{\mathit{cv}} = 0$;
    output to $\mathit{cp}(f,\mathit{sid}_n')$ weight 1.
    Annotated as SignalLoss detection point. \\[4pt]

% ── Condvar notify_all ───────────────────────────────
$(\mathit{sid}_n,\; \textsf{notify\_all}(\mathit{cv}),\;
  \textsf{next}(\mathit{sid}_n'))$
  with wait sites $\{w_1, \ldots, w_k\}$
  & Two transitions.
    \newline
    $t_{\mathit{notifyAll}}$ [CondvarNotifyAll]:
    input from $\mathit{cp}(f,\mathit{sid}_n)$ weight 1;
    guard $\mathit{nw}_{\mathit{cv}} > 0$;
    output to $\mathit{cp}(f,\mathit{sid}_n')$ weight 1;
    update $\mathit{na}_{w_1} \leftarrow \mathsf{true},\;
    \ldots,\; \mathit{na}_{w_k} \leftarrow \mathsf{true}$.
    \newline
    $t_{\mathit{allLost}}$ [CondvarNotifyAllLost]:
    input from $\mathit{cp}(f,\mathit{sid}_n)$ weight 1;
    guard $\mathit{nw}_{\mathit{cv}} = 0$;
    output to $\mathit{cp}(f,\mathit{sid}_n')$ weight 1.
    Annotated as SignalLoss detection point. \\[4pt]

% ── Spawn / Join ─────────────────────────────────────
$(\mathit{sid},\; \textsf{spawn}(g),\;
  \textsf{next}(\mathit{sid}'))$
  & One transition $t$ [Spawn]:
    input from $\mathit{cp}(f,\mathit{sid})$ weight 1;
    output to $\mathit{cp}(f,\mathit{sid}')$ weight 1 and
    $\mathit{cp}(g,\mathit{sid}_0^g)$ weight 1. \\[4pt]

$(\mathit{sid},\; \textsf{join}(g),\;
  \textsf{next}(\mathit{sid}'))$
  & One transition $t$ [Join]:
    input from $\mathit{cp}(f,\mathit{sid})$ weight 1 and
    $\mathit{cp}(g,\mathit{ret})$ weight 1;
    output to $\mathit{cp}(f,\mathit{sid}')$ weight 1. \\[4pt]

% ── Call ─────────────────────────────────────────────
$(\mathit{sid},\; \textsf{call}(g),\;
  \textsf{next}(\mathit{sid}'))$
  where $g$ has body
  & Call transition: $\mathit{cp}(f,\mathit{sid}) \to
    \mathit{cp}(g,\mathit{sid}_0^g)$.
    Return transition: $\mathit{cp}(g,\mathit{ret}) \to
    \mathit{cp}(f,\mathit{sid}')$. \\[4pt]

$(\mathit{sid},\; \textsf{call}(g),\;
  \textsf{next}(\mathit{sid}'))$
  where $g$ has summary only
  & One transition $t$ [Call]:
    input from $\mathit{cp}(f,\mathit{sid})$ weight 1;
    output to $\mathit{cp}(f,\mathit{sid}')$ weight 1.
    For each $x \in \Sigma(g).\mathit{writes}$:
    update $V[x] \leftarrow \top$. \\[4pt]

% ── Return ───────────────────────────────────────────
$(\mathit{sid},\; \textsf{return},\; \textsf{done})$
  & One transition $t$:
    input from $\mathit{cp}(f,\mathit{sid})$ weight 1;
    output to $\mathit{cp}(f,\mathit{ret})$ weight 1. \\[4pt]

% ── Switch ───────────────────────────────────────────
$(\mathit{sid},\; \textsf{read}(x),\;
  \textsf{switch}(x, [(v_1,s_1), \ldots, (v_k,s_k)]))$
  & $k$ transitions sharing input from $\mathit{cp}(f,\mathit{sid})$.
    Transition $t_i$ [Switch]: guard $V[x] = v_i$, output to
    $\mathit{cp}(f,s_i)$. \\

\bottomrule
\end{tabular}
\end{table*}

\subsubsection{Design Constraints}

The translation obeys four constraints.

\begin{enumerate}[leftmargin=1.4em]
\item \textbf{Faithful 1:1 mapping.} Every \textsc{Cir} statement produces at least one \textsc{Cvn} transition. No statements are merged, eliminated, or optimized.

\item \textbf{Protection is not translated.} The protection mapping $G$ is used only for \textsc{Cir} static checking. It does not produce any places, transitions, or arcs in the \textsc{Cvn}.

\item \textbf{Mode is not translated.} The \textsf{Sync}/\textsf{Async} mode of a resource does not affect the \textsc{Cvn} structure. Both modes produce the same places and transitions.

\item \textbf{Read with sequential transfer preserves anchors.} A $\textsf{read}$ followed by $\textsf{next}$ generates a pass-through transition with no guard and no update. This transition exists solely to maintain the anchor mapping: every \textsc{Cir} statement must correspond to at least one \textsc{Cvn} transition so that counterexamples can reference the statement.
\end{enumerate}

\subsubsection{Translation Error Taxonomy}

Table~\ref{tab:trans-errors} lists the translation error categories with prefix T.

\begin{table}[t]
\caption{Translation error categories.}
\label{tab:trans-errors}
\centering\footnotesize
\begin{tabular}{@{}llp{3.2cm}@{}}
\toprule
\textbf{Code} & \textbf{Category} & \textbf{Examples} \\
\midrule
T0xx & Input errors
  & \textsc{Cir} fails static check \\
T1xx & Resource errors
  & Unknown resource kind \\
T2xx & Control flow
  & Dangling transfer target \\
T3xx & Consistency
  & Incomplete anchor mapping \\
\bottomrule
\end{tabular}
\end{table}

\subsubsection{Translation Example}

Figure~\ref{fig:translation-example} illustrates the translation of the \texttt{notifier} function from the running example in Figure~\ref{fig:running-cir}. The function has five statements. The translation produces five control places, four transitions, and arcs connecting them. The \textsf{lock} transition at $\mathit{sid}$ \texttt{n1} consumes a token from $\mathit{rp}(\texttt{m0})$, modeling mutex acquisition. The \textsf{drop} transition at $\mathit{sid}$ \texttt{n4} returns the token, modeling mutex release. The \textsf{write} transition at $\mathit{sid}$ \texttt{n3} carries an update $V[\texttt{ready}] \leftarrow \mathsf{true}$.

\begin{figure}[t]
\centering
\small
\begin{verbatim}
Places:
  cp(notifier, n1)  cp(notifier, n2)
  cp(notifier, n3)  cp(notifier, n4)
  cp(notifier, n5)  cp(notifier, ret)
  rp(m0)  [shared with other functions]

Transitions:
  t_n1 [Lock]:
    in:  cp(notifier,n1) x1, rp(m0) x1
    out: cp(notifier,n2) x1
    anchor: {n1}

  t_n2 [CondvarNotify]:
    in:  cp(notifier,n2) x1, wp(w3) x1
    out: cp(notifier,n3) x1
    anchor: {n2}
    (Fires only if a waiter exists at w3.
     If wp(w3) is empty, notification is lost
     and a separate no-waiter transition
     advances control.)

  t_n3 [VarWrite]:
    in:  cp(notifier,n3) x1
    out: cp(notifier,n4) x1
    update: V[ready] <- true
    anchor: {n3}

  t_n4 [Unlock]:
    in:  cp(notifier,n4) x1
    out: cp(notifier,n5) x1, rp(m0) x1
    anchor: {n4}

  t_n5 [Return]:
    in:  cp(notifier,n5) x1
    out: cp(notifier,ret) x1
    anchor: {n5}
\end{verbatim}
\caption{Translation of the \texttt{notifier} function.}
\label{fig:translation-example}
\end{figure}



\section{Analysis and Repair}
\label{sec:analysis-repair}

The previous section defined the translation from \textsc{Cir} to \textsc{Cvn}. This section describes how the \textsc{Cvn} is analyzed and how analysis results are transformed into actionable repair guidance for the LLM. The analysis algorithms themselves are standard. The contribution lies in the connection between \textsc{Cvn} states and \textsc{Cir}-anchored counterexamples, and in the structured repair loop that exploits this connection.

% ─────────────────────────────────────────────────────────────
\subsection{State-Space Search}
\label{subsec:search}

The analysis engine explores the reachable state space of the \textsc{Cvn} starting from the initial state $(I_m, I_v)$. A state is a pair $(M, V)$ of marking and variable store. The engine
supports two traversal strategies: breadth-first search (BFS), which produces shortest counterexamples, and depth-first search (DFS), which uses less memory. Three termination parameters are configurable: a maximum number of visited states, a wall-clock timeout, and a flag indicating whether to stop after the first bug or to continue searching for all bugs.

At each state, the engine computes the set of enabled transitions (Definition~\ref{def:enabling}), fires each enabled transition (Definition~\ref{def:firing}), and checks whether the successor state has been visited. Visited states are identified by equality of both the marking $M$ and the variable store $V$ (Definition~\ref{def:state}). The state space is finite (Theorem~\ref{thm:finite}), so the search terminates.

Algorithm~\ref{alg:search} summarizes the procedure. Bug detection is interleaved with exploration: each newly reached state is checked against the bug patterns defined in the next subsection.

\begin{algorithm}[t]
\caption{State-space search with bug detection.}
\label{alg:search}
\begin{algorithmic}[1]
\Require \textsc{Cvn} $\mathcal{N}$ with initial state $(I_m, I_v)$
\Ensure Set of structured counterexamples, or ``verified''
\State $Q \gets [(I_m, I_v)]$
\State $\mathit{Visited} \gets \{(I_m, I_v)\}$
\State $\mathit{Parent} \gets \emptyset$
  \Comment{Records predecessor and transition for trace reconstruction}
\State $\mathit{Bugs} \gets \emptyset$
\While{$Q \neq \emptyset$ \textbf{and} within limits}
  \State $(M, V) \gets \mathsf{dequeue}(Q)$
  \State $E \gets \{t \in T \mid t \text{ is enabled at } (M, V)\}$
  \If{$\neg\,\mathsf{terminal}(M) \;\wedge\; E = \emptyset$}
    \State $\mathit{Bugs} \gets \mathit{Bugs} \cup
      \{\mathsf{build\_cex}(\textsc{Deadlock}, M, V, \mathit{Parent})\}$
    \If{stop-on-first} \textbf{break} \EndIf
  \EndIf
  \ForAll{$t \in E$}
    \State $(M', V') \gets \mathsf{fire}(t, M, V)$
    \If{$(M', V') \notin \mathit{Visited}$}
      \State $\mathit{Visited} \gets \mathit{Visited} \cup \{(M', V')\}$
      \State $\mathit{Parent}[(M', V')] \gets ((M, V),\; t)$
      \State $\mathsf{enqueue}(Q, (M', V'))$
    \EndIf
  \EndFor
\EndWhile
\State $\mathit{Bugs} \gets \mathit{Bugs} \cup
  \mathsf{scc\_liveness}(\mathit{Visited}, \mathit{Parent})$
\State \Return $\mathit{Bugs}$ if nonempty, else ``verified''
\end{algorithmic}
\end{algorithm}

% ─────────────────────────────────────────────────────────────
\subsection{Bug Detection Patterns}
\label{subsec:bug-patterns}

The engine detects five categories of concurrency bugs. Each category is defined as a predicate over states or state-graph structure.

\begin{definition}[Terminal state]
\label{def:terminal}
A state $(M, V)$ is terminal if every control token resides in a return place: for every function $f$ and every statement $\mathit{sid} \neq \mathit{ret}$, $M(\mathit{cp}(f, \mathit{sid})) = 0$.
\end{definition}

\begin{definition}[Deadlock]
\label{def:deadlock}
A state $(M, V)$ is a deadlock if it is not terminal and no transition is enabled:
\[
  \neg\,\mathsf{terminal}(M) \;\wedge\;
  \{t \in T \mid t \text{ is enabled at } (M,V)\} = \emptyset
\]
\end{definition}

\begin{definition}[Signal loss]
\label{def:signal-loss}
A signal loss occurs when a \textsf{CondvarNotify} or \textsf{CondvarNotifyAll} transition fires but no token exists in any wait place associated with that condition variable. The notification reaches no thread and is irrecoverably lost.
\end{definition}

The remaining three bug kinds are detected by analyzing the reachability graph after exploration completes.

\begin{definition}[Starvation]
\label{def:starvation}
A thread suffers starvation if there exists a nontrivial strongly connected component (SCC) in the reachability graph such that the thread's control token never advances within the SCC. The thread is alive (not in a return place) but makes no progress.
\end{definition}

\begin{definition}[Livelock]
\label{def:livelock}
A livelock exists if there is a nontrivial SCC in which all states are non-terminal. Threads continue to execute transitions but no thread ever reaches its return place.
\end{definition}

\begin{definition}[Channel blocking]
\label{def:channel-block}
A channel blocking deadlock is a special case of Definition~\ref{def:deadlock} in which the deadlock involves a \textsf{Recv} transition waiting for a channel token that will never
be produced because the sending thread is blocked on another resource.
\end{definition}

Signal loss is detected during exploration (line 12 in Algorithm~\ref{alg:search} can be extended to check the notification pattern upon firing). Starvation and livelock are detected by a post-exploration SCC analysis using Tarjan's algorithm~\cite{Tarjan1972} on the visited state graph.

Table~\ref{tab:bug-detection} summarizes the detection method for each bug kind.

\begin{table}[t]
\caption{Bug detection methods.}
\label{tab:bug-detection}
\centering\footnotesize
\begin{tabular}{@{}lll@{}}
\toprule
\textbf{Bug kind} & \textbf{Property} & \textbf{Detection} \\
\midrule
Deadlock       & Safety   & During exploration \\
SignalLoss     & Safety   & During exploration \\
ChannelBlock   & Safety   & During exploration \\
Starvation     & Liveness & Post-exploration SCC \\
Livelock       & Liveness & Post-exploration SCC \\
\bottomrule
\end{tabular}
\end{table}

% ─────────────────────────────────────────────────────────────
\subsection{Structured Counterexample}
\label{subsec:counterexample}

When a bug is detected, the engine constructs a structured counterexample. The counterexample is the primary interface between the formal verification engine and the LLM. Its design is driven by
two requirements: it must be precise enough for targeted repair, and it must be interpretable by an LLM without knowledge of \textsc{Cvn} internals.

\begin{definition}[Structured counterexample]
\label{def:counterexample}
A structured counterexample is a six-tuple
\[
  \mathit{CEx} = (\kappa,\; \pi,\;
    (M_f, V_f),\; R_{\mathit{inv}},\;
    F_{\mathit{inv}},\; S)
\]
where
\begin{itemize}[leftmargin=1.4em]
  \item $\kappa \in \{\textsc{Deadlock},\; \textsc{SignalLoss},\; \textsc{Starvation},\; \textsc{Livelock},\; \textsc{ChannelBlock}\}$ is the bug kind,
  \item $\pi = [(t_1, \mathit{sid}_1),\; \ldots,\; (t_n, \mathit{sid}_n)]$ is the transition trace, where each step carries the anchor mapping to \textsc{Cir} statement identifiers,
  \item $(M_f, V_f)$ is the final state at which the bug is detected,
  \item $R_{\mathit{inv}}$ is the set of resources involved in the bug,
  \item $F_{\mathit{inv}}$ is the set of functions involved in the bug, and
  \item $S$ is a template repair suggestion.
\end{itemize}
\end{definition}

The trace $\pi$ is reconstructed by following the $\mathit{Parent}$ map from the bug state back to the initial state and reversing the path. For liveness bugs detected by SCC analysis, the trace is a cycle within the SCC.

\paragraph{Trace anchoring.}
Each transition $t_i$ in the trace has an anchor $\mu(t_i) \in \mathit{SID}^+$ that maps it to one or more \textsc{Cir} statement identifiers. The LLM sees \textsc{Cir} identifiers, not \textsc{Cvn} transition names. This anchoring is maintained by the faithful 1:1 translation
(Section~\ref{subsec:translation}): every \textsc{Cvn} transition corresponds to exactly one \textsc{Cir} statement.

\paragraph{Final state interpretation.}
The final state $(M_f, V_f)$ is interpreted to produce a human-readable description. For a deadlock, the interpretation lists each thread's position (which control place holds a token), which resources each thread holds (which resource places have reduced token counts), and which resources each thread is waiting for (which input arcs of the blocked transitions require unavailable tokens). For a signal loss, the interpretation identifies the notification
site and the empty wait places.

\paragraph{Template repair suggestions.}
Each bug kind is associated with a repair template (Table~\ref{tab:repair-templates}). The template provides structural guidance rather than a concrete code patch. It identifies the nature
of the fix and the affected \textsc{Cir} statements.

\begin{table}[t]
\caption{Repair suggestion templates.}
\label{tab:repair-templates}
\centering\footnotesize
\begin{tabular}{@{}lp{5cm}@{}}
\toprule
\textbf{Bug kind} & \textbf{Suggestion template} \\
\midrule
\textsc{Deadlock}
  & Unify the lock acquisition order across all involved
    functions. Suggested order: $r_1 < r_2 < \ldots$ \\[3pt]
\textsc{SignalLoss}
  & Ensure the predicate is established before the notification.
    Protect \textsf{wait} with a while-loop that re-checks the
    condition. \\[3pt]
\textsc{ChannelBlock}
  & Do not perform blocking channel operations while holding a
    lock. Release the lock before \textsf{send} or \textsf{recv}. \\[3pt]
\textsc{Starvation}
  & Reduce lock hold duration. Consider bounded retry or fairness
    mechanisms. \\[3pt]
\textsc{Livelock}
  & Introduce backoff or a deterministic coordination protocol. \\
\bottomrule
\end{tabular}
\end{table}

\paragraph{Counterexample format.}
The counterexample is serialized in two formats. A structured format (for programmatic consumption) contains all fields of Definition~\ref{def:counterexample} as typed data. A text format (for LLM consumption) presents the same information in a template:

\begin{lstlisting}
BUG: <kind>
TRACE:
  step <i>: [<function>/<sid>] <operation> <resource>
  ...
  >> <bug-specific annotation>
STATE:
  <thread positions and held resources>
RESOURCES: <involved resources>
FUNCTIONS: <involved functions>
SUGGESTION:
  <template text with affected sids>
\end{lstlisting}

The text format uses \textsc{Cir} vocabulary throughout. The LLM does not need to know what a \textsc{Cvn} place or transition is. It sees function names, statement identifiers, operation names, and resource names, all of which match the \textsc{Cir} it generated.

% ─────────────────────────────────────────────────────────────
\subsection{Repair Loop}
\label{subsec:repair-loop}

The repair loop connects the verification engine back to the LLM, closing the generate-verify-repair cycle described in Section~\ref{sec:architecture}.

\subsubsection{Repair Strategy}

Two classes of errors require different repair strategies.

\textbf{Layer-1 errors} (CIR static check, E0xx--E8xx) are structural. Some admit local automatic repair. A missing \textsf{drop} statement can be inserted after the last use of the
lock guard. A duplicate resource name can be renamed. These repairs do not require semantic understanding.

\textbf{Layer-2 errors} (CVN model checking) are behavioral. They arise from cross-thread interactions and require semantic understanding to repair. Fixing a deadlock requires choosing between several strategies: unifying lock order, reducing lock granularity, or restructuring the synchronization protocol. This choice depends on the program's business intent, which only the LLM has access to. Layer-2 errors are therefore always repaired by the LLM.

\subsubsection{Repair Prompt Construction}

When a Layer-2 bug is found, the engine constructs a repair prompt from the structured counterexample. The prompt contains five sections.

\begin{enumerate}[leftmargin=1.4em]
\item \textbf{Context.} The original \textsc{Cir} artifact that was verified, so the LLM can see the full concurrency structure.

\item \textbf{Bug report.} The text-format counterexample from Section~\ref{subsec:counterexample}, including the bug kind, the trace with statement identifiers, and the final state.

\item \textbf{Repair suggestion.} The template suggestion from Table~\ref{tab:repair-templates}, instantiated with the specific resources and statement identifiers from the counterexample.

\item \textbf{Constraints.} Instructions to the LLM: minimize edits, preserve the resource naming, do not introduce new synchronization primitives unless the suggestion requires it, and maintain statement identifier stability where possible.

\item \textbf{Output format.} A specification of the expected \textsc{Cir} output format, so the repaired artifact can be parsed and re-verified.
\end{enumerate}

\subsubsection{Iterative Repair}

Algorithm~\ref{alg:repair} describes the iterative repair process.

\begin{algorithm}[t]
\caption{Generate-verify-repair loop.}
\label{alg:repair}
\begin{algorithmic}[1]
\Require Source code or specification
\Ensure Verified \textsc{Cir} or failure
\State LLM generates \textsc{Cir} from source
\For{$\mathit{round} = 1$ \textbf{to} $K$}
  \State Run \textsc{Cir} static checker
  \If{Layer-1 errors}
    \If{auto-repairable}
      \State Apply local fix
    \Else
      \State Send error report to LLM
      \State LLM regenerates \textsc{Cir}
    \EndIf
    \State \textbf{continue}
  \EndIf
  \State Translate \textsc{Cir} to \textsc{Cvn}
  \State Run state-space search
  \If{no bugs found}
    \State \Return verified \textsc{Cir}
  \EndIf
  \State Construct repair prompt from counterexample
  \State LLM generates repaired \textsc{Cir}
\EndFor
\State \Return failure after $K$ rounds
\end{algorithmic}
\end{algorithm}

The maximum number of rounds $K$ is a configurable parameter. Each round involves one LLM invocation and one verification pass. In practice, $K = 3$ to $5$ is sufficient for most patterns, because the structured counterexample provides precise enough guidance for the LLM to produce a correct fix within a small number of attempts.

\subsubsection{Convergence Considerations}

The repair loop is not guaranteed to converge. The LLM may introduce new bugs while fixing existing ones. Two mechanisms mitigate this risk. First, the \textsc{Cir} static checker (Layer~1) catches structural regressions immediately. A repair that introduces an undefined
resource reference or a broken control-flow edge is rejected before reaching the more expensive \textsc{Cvn} analysis. Second, the repair prompt explicitly instructs the LLM to minimize
edits. This reduces the probability of unintended side effects. The prompt also includes the full \textsc{Cir} as context, so the LLM can reason about the impact of its changes on the rest of the concurrency structure. If the loop does not converge within $K$ rounds, the system reports
all discovered bugs and the history of repair attempts, allowing a human developer to intervene.




% ══════════════════════════════════════════════════════════════
%  §6  Formal Properties and Discussion
% ══════════════════════════════════════════════════════════════
\section{Formal Properties and Discussion}
\label{sec:properties}

This section establishes the formal relationship between \textsc{Cir}
executions and \textsc{Cvn} firings, states the guarantees provided
by the analysis, and discusses the modeling decisions and limitations
of the approach.

% ─────────────────────────────────────────────────────────────
\subsection{CIR Execution Semantics}
\label{subsec:cir-semantics}

\begin{definition}[\textsc{Cir} configuration]
\label{def:cir-config}
A \textsc{Cir} configuration is a tuple
$C = (\mathit{pc}, L, W, \mathit{NW}, \mathit{NA}, V)$
where
\begin{itemize}[leftmargin=1.4em]
  \item $\mathit{pc} : \mathit{ThreadId} \to \mathit{SID} \cup
        \{\mathit{ret}\}$ maps each thread to its current statement,
  \item $L : \mathit{ResourceName} \to \mathbb{N}_0$ records the
        available count of each synchronization resource (mutexes,
        semaphores, channels) and each condvar notification token,
  \item $W : \mathit{SID} \to \mathbb{N}_0$ records the number of
        threads blocked at each wait site,
  \item $\mathit{NW} : \mathit{CondvarName} \to \mathbb{N}_0$
        records the total number of waiters for each condvar,
  \item $\mathit{NA} : \mathit{SID} \to \mathsf{Bool}$ records the
        notify-all flag for each wait site, and
  \item $V : \mathit{VarName} \to \mathit{Val}$ is the variable
        store.
\end{itemize}
\end{definition}

The \textsc{Cir} step rules for condvar operations are:

\textbf{Wait.} Thread $\theta$ at $(\mathit{sid},\;
\textsf{wait}(\mathit{cv}, m),\; \textsf{next}(\mathit{sid}'))$
performs two sub-steps. Enter: $L'(m) = L(m) + 1$,
$W'(\mathit{sid}) = W(\mathit{sid}) + 1$,
$\mathit{NW}'(\mathit{cv}) = \mathit{NW}(\mathit{cv}) + 1$,
$\mathit{NA}'(\mathit{sid}) = \mathsf{false}$,
thread enters waiting state at $\mathit{sid}$. Wake (by notify):
requires $L(\mathit{cv}) \geq 1$; sets $L'(\mathit{cv}) = L(\mathit{cv}) - 1$,
$W'(\mathit{sid}) = W(\mathit{sid}) - 1$,
$\mathit{NW}'(\mathit{cv}) = \mathit{NW}(\mathit{cv}) - 1$,
thread enters reacquire state. Wake (by notify-all):
requires $\mathit{NA}(\mathit{sid}) = \mathsf{true}$; sets
$W'(\mathit{sid}) = W(\mathit{sid}) - 1$,
$\mathit{NW}'(\mathit{cv}) = \mathit{NW}(\mathit{cv}) - 1$,
$\mathit{NA}'(\mathit{sid}) = \mathsf{false}$,
thread enters reacquire state. Reacquire: requires $L(m) \geq 1$;
sets $L'(m) = L(m) - 1$, $\mathit{pc}'(\theta) = \mathit{sid}'$.

\textbf{Notify\_one.} If $\mathit{NW}(\mathit{cv}) > 0$: set
$L'(\mathit{cv}) = L(\mathit{cv}) + 1$,
$\mathit{pc}'(\theta) = \mathit{sid}'$. If
$\mathit{NW}(\mathit{cv}) = 0$: notification lost,
$\mathit{pc}'(\theta) = \mathit{sid}'$.

\textbf{Notify\_all.} If $\mathit{NW}(\mathit{cv}) > 0$: for each
wait site $w_i$ of $\mathit{cv}$, set
$\mathit{NA}'(w_i) = \mathsf{true}$,
$\mathit{pc}'(\theta) = \mathit{sid}'$. If
$\mathit{NW}(\mathit{cv}) = 0$: notification lost,
$\mathit{pc}'(\theta) = \mathit{sid}'$.


% ─────────────────────────────────────────────────────────────
\subsection{Encoding Function}
\label{subsec:encoding}

\begin{definition}[Encoding]
\label{def:encoding}
The encoding function $\mathsf{enc}$ maps a \textsc{Cir}
configuration $C = (\mathit{pc}, L, W, \mathit{NW}, \mathit{NA}, V)$
to a \textsc{Cvn} state $(M, V')$ as follows:
\begin{itemize}[leftmargin=1.4em]
  \item For each thread $\theta$ at $\mathit{sid}$ in function $f$
        (not waiting, not reacquiring):
        $M(\mathit{cp}(f, \mathit{sid})) \mathrel{+}= 1$.
  \item For each thread $\theta$ waiting at site $\mathit{sid}$:
        $M(\mathit{wp}(\mathit{sid})) \mathrel{+}= 1$.
  \item For each thread $\theta$ reacquiring at site $\mathit{sid}$:
        $M(\mathit{ra}(\mathit{sid})) \mathrel{+}= 1$.
  \item For each synchronization resource $r$:
        $M(\mathit{rp}(r)) = L(r)$.
  \item For each condvar $\mathit{cv}$:
        $M(\mathit{rp}(\mathit{cv})) = L(\mathit{cv})$.
  \item $V'$ extends $V$ with:
        $V'[\mathit{nw}_{\mathit{cv}}] = \mathit{NW}(\mathit{cv})$
        for each condvar, and
        $V'[\mathit{na}_{w_i}] = \mathit{NA}(w_i)$ for each wait
        site.
\end{itemize}
\end{definition}

The encoding preserves all concurrency-relevant information. It does
not preserve thread identity (which thread is at which position),
because \textsc{Cvn} tokens are indistinguishable. This is
acceptable because the analysis does not require thread identity:
two tokens in the same control place represent two interchangeable
threads at the same program point.

% ─────────────────────────────────────────────────────────────

\subsection{Step Correspondence}
\label{subsec:step-correspondence}

\begin{theorem}[Forward simulation]
\label{thm:forward}
Let $C$ be a \textsc{Cir} configuration and
$(M, V) = \mathsf{enc}(C)$. For every \textsc{Cir} step
$C \xrightarrow{o} C'$, there exists a \textsc{Cvn} transition $t$
such that $(M, V) \xrightarrow{t} (M', V')$ and
$(M', V') = \mathsf{enc}(C')$.
\end{theorem}

\begin{proof}
By case analysis on the operation $o$. We present the condvar cases;
the remaining cases follow the same structure as before.

\textbf{Case} $o = \textsf{lock}(m)$. The \textsc{Cir} step
requires $L(m) \geq 1$ and produces $L'(m) = L(m) - 1$,
$\mathit{pc}'(\theta) = \mathit{sid}'$. The corresponding
\textsc{Cvn} transition $t$ [Lock] has input arcs from
$\mathit{cp}(f, \mathit{sid})$ weight 1 and $\mathit{rp}(m)$
weight 1, and output arc to $\mathit{cp}(f, \mathit{sid}')$
weight 1. The precondition $M(\mathit{rp}(m)) = L(m) \geq 1$
guarantees that $t$ is enabled. Firing $t$ produces
$M'(\mathit{rp}(m)) = M(\mathit{rp}(m)) - 1 = L'(m)$ and
$M'(\mathit{cp}(f, \mathit{sid}')) = M(\mathit{cp}(f,
\mathit{sid}')) + 1$, matching $\mathsf{enc}(C')$.

\textbf{Case} $o = \textsf{drop}(m)$. Symmetric: the transition
returns a token to $\mathit{rp}(m)$.

\textbf{Case} $o = \textsf{write}(x, v)$. The \textsc{Cir} step
sets $V'[x] = v$. The \textsc{Cvn} transition [VarWrite] carries
the update $V[x] \leftarrow v$, producing the same variable
store.

\textbf{Case} $o = \textsf{read}(x)$ with
$\textsf{branch}(\beta, \mathit{sid}_t, \mathit{sid}_f)$. The
\textsc{Cir} step evaluates $\beta$ to a definite boolean and
branches accordingly. The \textsc{Cvn} has two transitions
$t_\top$ and $t_\bot$ with complementary guards. Since $\beta$
evaluates to a concrete boolean, exactly one guard evaluates to
$\mathsf{true}$ and the other to $\mathsf{false}$. The enabled
transition matches the \textsc{Cir} branch direction.

\textbf{Case} $o = \textsf{wait}(\mathit{cv}, m)$. The
\textsc{Cir} step moves the thread to a waiting state and releases
the mutex. The \textsc{Cvn} transition $t_w$ [CondvarWait] moves
the token from $\mathit{cp}(f, \mathit{sid})$ to
$\mathit{wp}(\mathit{sid})$ and returns a token to
$\mathit{rp}(m)$, matching the encoding of the new configuration.

\textbf{Case} $o = \textsf{notify\_one}(\mathit{cv})$. If at
least one thread is waiting (some $W(\mathit{sid}_w) \geq 1$),
the \textsc{Cir} step wakes one waiter nondeterministically. The
corresponding wake transition $t_i$ [CondvarNotify] consumes a
token from $\mathit{wp}(\mathit{sid}_w)$, matching the waiter
selection. If no thread is waiting, the \textsc{Cir} step has no
wake effect and advances control. The corresponding pass-through
transition advances control without consuming from any wait place.

\textbf{Case} $o = \textsf{spawn}(g)$. The \textsc{Cir} step
creates a new thread at the entry of $g$. The \textsc{Cvn}
transition [Spawn] produces tokens in both
$\mathit{cp}(f, \mathit{sid}')$ and
$\mathit{cp}(g, \mathit{sid}_0^g)$, matching the fork.

\textbf{Case} $o = \textsf{wait}(\mathit{cv}, m)$, \emph{enter
sub-step}. The \textsc{Cir} step increments
$W(\mathit{sid})$, increments $\mathit{NW}(\mathit{cv})$,
resets $\mathit{NA}(\mathit{sid})$, and releases the mutex
($L(m) + 1$). The \textsc{Cvn} transition
$t_{\mathit{enter}}$ [CondvarWaitEnter] consumes one token from
$\mathit{cp}(f, \mathit{sid})$, produces one token in
$\mathit{wp}(\mathit{sid})$ and one token in $\mathit{rp}(m)$,
updates $V[\mathit{nw}_{\mathit{cv}}] \leftarrow
V[\mathit{nw}_{\mathit{cv}}] + 1$ and
$V[\mathit{na}_{\mathit{sid}}] \leftarrow \mathsf{false}$.
Enabling: $M(\mathit{cp}(f, \mathit{sid})) \geq 1$ holds because
the thread is at $\mathit{sid}$. No guard. The resulting marking and
variable store match $\mathsf{enc}(C')$.

\textbf{Case} $o = \textsf{wait}(\mathit{cv}, m)$, \emph{wake by
notify sub-step}. The \textsc{Cir} step requires
$L(\mathit{cv}) \geq 1$, decrements $W(\mathit{sid})$, decrements
$\mathit{NW}(\mathit{cv})$, and consumes one notification token. The
\textsc{Cvn} transition $t_{\mathit{wake1}}$
[CondvarWakeByNotify] consumes one token from
$\mathit{wp}(\mathit{sid})$ and one token from
$\mathit{rp}(\mathit{cv})$, produces one token in
$\mathit{ra}(\mathit{sid})$, updates
$V[\mathit{nw}_{\mathit{cv}}] \leftarrow
V[\mathit{nw}_{\mathit{cv}}] - 1$. Enabling:
$M(\mathit{wp}(\mathit{sid})) \geq 1$ holds because the thread is
waiting, and $M(\mathit{rp}(\mathit{cv})) = L(\mathit{cv}) \geq 1$.
No guard. Match follows.

\textbf{Case} $o = \textsf{wait}(\mathit{cv}, m)$, \emph{wake by
notify-all sub-step}. The \textsc{Cir} step requires
$\mathit{NA}(\mathit{sid}) = \mathsf{true}$, does not consume a
notification token, decrements $W(\mathit{sid})$ and
$\mathit{NW}(\mathit{cv})$, resets
$\mathit{NA}(\mathit{sid})$. The \textsc{Cvn} transition
$t_{\mathit{wakeA}}$ [CondvarWakeByNotifyAll] consumes one token
from $\mathit{wp}(\mathit{sid})$, has guard
$V[\mathit{na}_{\mathit{sid}}] = \mathsf{true}$, produces one
token in $\mathit{ra}(\mathit{sid})$, and updates
$V[\mathit{nw}_{\mathit{cv}}] \leftarrow
V[\mathit{nw}_{\mathit{cv}}] - 1$,
$V[\mathit{na}_{\mathit{sid}}] \leftarrow \mathsf{false}$. Enabling:
token in $\mathit{wp}(\mathit{sid})$ holds; guard holds because
$V[\mathit{na}_{\mathit{sid}}] = \mathit{NA}(\mathit{sid}) =
\mathsf{true}$. Match follows.

\textbf{Case} $o = \textsf{wait}(\mathit{cv}, m)$, \emph{reacquire
sub-step}. The \textsc{Cir} step requires $L(m) \geq 1$ and sets
$\mathit{pc}'(\theta) = \mathit{sid}'$. The \textsc{Cvn} transition
$t_{\mathit{reacq}}$ [CondvarReacquire] consumes one token from
$\mathit{ra}(\mathit{sid})$ and one from $\mathit{rp}(m)$, produces
one token in $\mathit{cp}(f, \mathit{sid}')$. Enabling: both tokens
are present. Match follows.

\textbf{Case} $o = \textsf{notify\_one}(\mathit{cv})$ with
$\mathit{NW}(\mathit{cv}) > 0$. The \textsc{Cir} step deposits one
notification token: $L'(\mathit{cv}) = L(\mathit{cv}) + 1$. The
\textsc{Cvn} transition $t_{\mathit{notify}}$ [CondvarNotify] has
guard $V[\mathit{nw}_{\mathit{cv}}] > 0$, which holds because
$V[\mathit{nw}_{\mathit{cv}}] = \mathit{NW}(\mathit{cv}) > 0$.
The transition produces one token in $\mathit{rp}(\mathit{cv})$.
Match follows.

\textbf{Case} $o = \textsf{notify\_one}(\mathit{cv})$ with
$\mathit{NW}(\mathit{cv}) = 0$. The notification is lost. The
\textsc{Cvn} transition $t_{\mathit{lost}}$ [CondvarNotifyLost]
has guard $V[\mathit{nw}_{\mathit{cv}}] = 0$, which holds. The
transition advances control without depositing a token. Match
follows.

\textbf{Case} $o = \textsf{notify\_all}(\mathit{cv})$ with
$\mathit{NW}(\mathit{cv}) > 0$. The \textsc{Cir} step sets
$\mathit{NA}'(w_i) = \mathsf{true}$ for all wait sites $w_i$ of
$\mathit{cv}$. The \textsc{Cvn} transition
$t_{\mathit{notifyAll}}$ [CondvarNotifyAll] has guard
$V[\mathit{nw}_{\mathit{cv}}] > 0$, which holds. The transition
updates $V[\mathit{na}_{w_i}] \leftarrow \mathsf{true}$ for all
$w_i$. Match follows.

\textbf{Case} $o = \textsf{notify\_all}(\mathit{cv})$ with
$\mathit{NW}(\mathit{cv}) = 0$. Same as the notify-lost case.

The remaining cases (\textsf{join}, \textsf{cas}, \textsf{acquire},
\textsf{release}, \textsf{send}, \textsf{recv}, \textsf{call},
\textsf{return}) follow the same pattern. Each \textsc{Cir}
precondition maps to the corresponding \textsc{Cvn} enabling
condition, and each postcondition maps to the corresponding firing
effect.
\end{proof}

The forward simulation guarantees that every \textsc{Cir} behavior
is represented in the \textsc{Cvn}. The \textsc{Cvn} cannot miss a
bug that exists in the \textsc{Cir}.

\begin{theorem}[Backward simulation under concrete valuation]
\label{thm:backward}
Let $(M, V) = \mathsf{enc}(C)$ where all variables in $V$ have
concrete values (no variable equals $\top$). For every enabled
\textsc{Cvn} transition $t$ at $(M, V)$, there exists a \textsc{Cir}
step $C \xrightarrow{o} C'$ such that
$\mathsf{enc}(C') = \mathsf{fire}(t, M, V)$.
\end{theorem}

\begin{proof}
When all variables are concrete, guard evaluation in the \textsc{Cvn}
produces only $\mathsf{true}$ or $\mathsf{false}$, never
$\mathsf{unknown}$. The enabling conditions of \textsc{Cvn}
transitions therefore coincide exactly with the preconditions of
\textsc{Cir} operations. The argument proceeds by case analysis,
symmetric to Theorem~\ref{thm:forward}.
\end{proof}

When some variable has the value $\top$, the backward simulation may
fail: a guard involving that variable evaluates to $\mathsf{unknown}$,
enabling the transition in the \textsc{Cvn}, but the corresponding
\textsc{Cir} branch may not be takeable under any concrete valuation.
This is the source of potential false positives, discussed in
Section~\ref{subsec:over-approx}.
% ─────────────────────────────────────────────────────────────
\subsection{Deadlock Correspondence}
\label{subsec:deadlock-correspondence}

\begin{theorem}[Deadlock correspondence]
\label{thm:deadlock-corr}
A \textsc{Cir} configuration $C$ is a deadlock (non-terminal with no
enabled step) if and only if $\mathsf{enc}(C)$ is a deadlock in the
\textsc{Cvn} (non-terminal marking with no enabled transition),
provided all variables have concrete values.
\end{theorem}

\begin{proof}
A configuration is a deadlock when the set of enabled
operations/transitions is empty and at least one thread has not
returned. By Theorems~\ref{thm:forward} and~\ref{thm:backward}, the
enabled sets coincide under concrete valuation. The non-terminal
condition is preserved by the encoding: a thread at
$\mathit{sid} \neq \mathit{ret}$ maps to a token in a non-return
control place.
\end{proof}

Under $\top$-valued variables, the correspondence is one-directional:
a \textsc{Cir} deadlock implies a \textsc{Cvn} deadlock (by
Theorem~\ref{thm:forward}), but a \textsc{Cvn} deadlock may not
correspond to any \textsc{Cir} deadlock (the \textsc{Cvn} may have
reached the state through an infeasible branch enabled by
$\mathsf{unknown}$).


\subsection{Signal Loss Soundness}
\label{subsec:signal-loss-soundness}

The updated condvar model enables a precise characterization of
signal loss.

\begin{proposition}[Signal loss detection]
\label{prop:signal-loss}
A \textsf{notify\_one} or \textsf{notify\_all} operation at
$\mathit{sid}_n$ causes a signal loss in a \textsc{Cir} execution if
and only if the corresponding \textsc{Cvn} transition
$t_{\mathit{lost}}$ or $t_{\mathit{allLost}}$ fires at state
$(M, V)$ with $V[\mathit{nw}_{\mathit{cv}}] = 0$.
\end{proposition}

\begin{proof}
Forward direction: if the \textsc{Cir} notification occurs with
$\mathit{NW}(\mathit{cv}) = 0$, then by the encoding
$V[\mathit{nw}_{\mathit{cv}}] = 0$, and the guard of
$t_{\mathit{lost}}$ is satisfied. Backward direction: if
$t_{\mathit{lost}}$ fires, then
$V[\mathit{nw}_{\mathit{cv}}] = 0 = \mathit{NW}(\mathit{cv})$,
and the \textsc{Cir} notification has no waiters.
\end{proof}

This proposition establishes that the signal-loss annotation on the
lost transitions is both sound (no false negatives) and precise (no
false positives) with respect to the \textsc{Cir} semantics. A
signal loss is recorded if and only if the notification has no
waiters in the corresponding \textsc{Cir} execution.



% ─────────────────────────────────────────────────────────────
\subsection{Soundness of Over-Approximation}
\label{subsec:over-approx}

\begin{theorem}[Soundness]
\label{thm:soundness}
If a \textsc{Cir} execution can reach a bug configuration (deadlock,
signal loss, starvation, or livelock), then the \textsc{Cvn}
state-space search will find a corresponding bug state.
\end{theorem}

\begin{proof}
Let $C_0 \xrightarrow{o_1} C_1 \xrightarrow{o_2} \cdots
\xrightarrow{o_n} C_n$ be a \textsc{Cir} execution reaching a bug
configuration $C_n$. By Theorem~\ref{thm:forward}, there exists a
corresponding \textsc{Cvn} firing sequence
$(M_0, V_0) \xrightarrow{t_1} \cdots \xrightarrow{t_n}
(M_n, V_n)$ with $(M_i, V_i) = \mathsf{enc}(C_i)$ for each $i$.

The $\mathsf{unknown}$ mechanism only adds enabled transitions to
the \textsc{Cvn} (both branches of an unknown guard become enabled).
It never removes transitions that would be enabled under concrete
evaluation. Therefore, every transition enabled in the concrete
\textsc{Cvn} is also enabled in the $\mathsf{unknown}$-augmented
\textsc{Cvn}. The firing sequence above remains valid.

The bug predicate (deadlock, signal loss, starvation, livelock) is
defined over \textsc{Cvn} states and state-graph structure
(Definitions~\ref{def:deadlock}--\ref{def:channel-block}). Since
$(M_n, V_n) = \mathsf{enc}(C_n)$ and $C_n$ satisfies the bug
predicate at the \textsc{Cir} level, $(M_n, V_n)$ satisfies the
corresponding predicate at the \textsc{Cvn} level.

The finite state space (Theorem~\ref{thm:finite}) guarantees that
the search terminates and visits $(M_n, V_n)$.
\end{proof}

The converse does not hold: the \textsc{Cvn} may report bugs in
states that are reachable only through $\mathsf{unknown}$-enabled
branches and that no concrete \textsc{Cir} execution can reach.
These are false positives. In practice, false positives arise only
when variables modified by summarized functions
(Section~\ref{subsec:translation}, Phase~3) are used in branch
conditions. For fully modeled programs where all functions have
bodies in the \textsc{Cir}, all variables have concrete values
throughout execution, and the analysis produces no false positives.

% ─────────────────────────────────────────────────────────────
\subsection{Loop Modeling}
\label{subsec:loop-discussion}

The system does not perform automated loop unrolling. This is a
deliberate design decision based on two observations.

\textbf{Concurrency bugs are topological.} Deadlocks, signal losses,
and starvation are triggered by the structure of synchronization
interactions, not by loop iteration counts. A circular lock
dependency causes deadlock on the first iteration. A misplaced
notification causes signal loss regardless of how many times the loop
body executes. Two or three concurrent instances are sufficient to
expose the vast majority of synchronization bugs.

\textbf{LLM-driven instance unfolding replaces loop unrolling.}
When the source code contains a loop that spawns $N$ workers, the
LLM unfolds the loop into a small number (typically 2 or 3) of
explicit \textsf{spawn} statements in the \textsc{Cir}. Each
instance receives a unique function name (e.g.,
\texttt{worker\_0}, \texttt{worker\_1}) and therefore a unique set
of control places in the \textsc{Cvn}. Black tokens suffice to
distinguish thread positions because the places themselves are
distinct.

Non-spawn loops retain their back-edge structure in \textsc{Cir}. A
loop body that modifies a shared variable may cause that variable to
take different values on different iterations. If the variable's
value cannot be determined statically, it becomes $\top$ in the
\textsc{Cvn} after the first iteration. The three-valued guard
evaluation then enables both branches of any condition involving
that variable, ensuring that all paths through the loop body are
explored. This provides a safe over-approximation of all loop
iterations without explicit unrolling.

\textbf{Consequence.} The system verifies a fixed, bounded number of
concurrent instances, not parametric correctness for arbitrary $N$.
Bugs that require more instances than the LLM generates (e.g., a
$k$-way circular deadlock where $k$ exceeds the unfolding count)
will not be detected. In practice, setting the unfolding count to
$\min(N, 3)$ is sufficient for the patterns in our test matrix
(Section~\ref{sec:evaluation}).

% ─────────────────────────────────────────────────────────────
\subsection{Why Colored Tokens Are Unnecessary}
\label{subsec:no-color}

Classical colored Petri nets (CPNs) attach typed data to tokens,
enabling a single parameterized net to model multiple instances. In
our architecture, this parameterization is unnecessary for two
reasons.

First, the LLM has already resolved all resource identities at the
\textsc{Cir} level. Tokens do not need to carry resource names
because each resource has a dedicated place. Second, the LLM has
already unfolded concurrent instances into distinct functions. Tokens
do not need to carry thread identities because each instance has
dedicated control places. Black tokens are sufficient.

This simplification eliminates the need for color sets, binding
expressions, arc inscriptions, and the canonical state
representation that CPNs require for state-space exploration. The
resulting \textsc{Cvn} is a standard weighted P/T net (augmented
with global-variable guards), for which efficient analysis algorithms
are well established~\cite{Murata1989}.

% ─────────────────────────────────────────────────────────────
\subsection{Limitations}
\label{subsec:limitations}

\textbf{LLM reliability.} The correctness of the \textsc{Cir}
depends on the LLM's understanding of the source code. Complex
aliasing through unsafe code, raw pointers, or foreign function
interfaces may exceed current LLM capabilities. Errors in the LLM's
resource identification manifest as structural errors in \textsc{Cir}
and are caught by the static checker, but subtle omissions (e.g.,
failing to model a shared resource) may lead to missed bugs.

\textbf{Over-approximation.} The $\top$ mechanism may produce false
positives when branch conditions involve variables whose values
cannot be determined. The false positive rate depends on the
proportion of summarized functions and the complexity of branch
conditions.

\textbf{Bounded instances.} The system verifies a fixed number of
unfolded instances. Parametric correctness (``safe for any $N$'') is
outside the current scope.

\textbf{No weak-memory model.} Atomic operations are modeled at the
variable level without a formal memory model. Bugs that depend on
weak-memory ordering (e.g., missing barriers on ARM) are not
detected.

\textbf{State explosion.} The state space grows exponentially with
the number of concurrent entities. For programs with many threads
and many synchronization operations, the search may exceed the
configured state limit before completing. The timeout and
state-limit parameters allow trading completeness for practicality.



% ══════════════════════════════════════════════════════════════
%  §7  Evaluation
% ══════════════════════════════════════════════════════════════
\section{Evaluation}
\label{sec:evaluation}

We evaluate the generate--verify--repair pipeline through four
research questions:

\begin{itemize}[leftmargin=1.4em]
  \item \textbf{RQ1 (CIR Generation).} Which LLMs can produce a
    syntactically valid \textsc{Cir} artifact, and how many
    iterative rounds of feedback are required?
  \item \textbf{RQ2 (Bug Detection \& Repair).} For a buggy
    \textsc{Cir}, can \textsc{Cvn} analysis detect the bug? What is
    the \textsc{Cvn} size and analysis cost? How many repair
    iterations does the LLM need, and does the repair process
    introduce regressions?
  \item \textbf{RQ3 (Translation Correctness).} Does the
    \textsc{Cir}$\,\to\,$\textsc{Cvn} translation preserve the
    structural invariants required by the step-correspondence
    theorem?
  \item \textbf{RQ4 (Code--CIR Faithfulness).} Can we guarantee
    that LLM-generated source code faithfully implements the
    verified \textsc{Cir}?
\end{itemize}

% ─────────────────────────────────────────────────────────────
\subsection{Test Matrix}
\label{subsec:test-matrix}

Table~\ref{tab:test-matrix} lists the 10 concurrency patterns used
across all RQs. Seven patterns contain known bugs and three are
bug-free baselines that verify the absence of false positives. Each
pattern exercises a distinct combination of synchronization
primitives and bug kinds.

\begin{table*}[t]
\caption{Test matrix. Each pattern is verified in both buggy and fixed
form. Detection layer indicates whether the bug is caught by
\textsc{Cir} static checking (Layer~1) or \textsc{Cvn} state-space
search (Layer~2).}
\label{tab:test-matrix}
\centering\footnotesize
\begin{tabular}{@{}clllllp{3.6cm}@{}}
\toprule
\textbf{\#} & \textbf{Pattern} & \textbf{Bug kind}
  & \textbf{Safety/Liveness} & \textbf{Layer} & \textbf{Key resources}
  & \textbf{Repair strategy} \\
\midrule
1  & Two-mutex deadlock
   & Deadlock & Safety & 2
   & Mutex $\times 2$
   & Unify lock order \\[2pt]
2  & Condvar signal loss
   & SignalLoss & Safety & 2
   & Condvar, Mutex, Var
   & Set predicate before notify \\[2pt]
3  & Channel + mutex deadlock
   & Deadlock & Safety & 2
   & Channel, Mutex
   & Release lock before channel op \\[2pt]
4  & Three-lock circular deadlock
   & Deadlock & Safety & 2
   & Mutex $\times 3$
   & Establish global lock order \\[2pt]
5  & Semaphore throttling
   & (none) & --- & ---
   & Semaphore
   & Baseline: no false positive \\[2pt]
6  & CAS contention
   & (none) & --- & ---
   & Atomic
   & Baseline: Unknown split \\[2pt]
7  & RwLock writer starvation
   & Starvation & Liveness & 2
   & RwLock, Var
   & Reduce read-lock hold time \\[2pt]
8  & Condvar spurious wakeup
   & SignalLoss & Safety & 2
   & Condvar, Mutex, Var
   & While-loop guard on wait \\[2pt]
9  & Dual condvar cross-wait
   & Deadlock & Safety & 2
   & Condvar $\times 2$, Mutex $\times 2$
   & Break wait cycle \\[2pt]
10 & FnSummary propagation
   & (none) & --- & ---
   & Var, Mutex
   & Baseline: Unknown coverage \\
\bottomrule
\end{tabular}
\end{table*}

% ─────────────────────────────────────────────────────────────
\subsection{Pattern Descriptions}
\label{subsec:pattern-descriptions}

\textbf{Pattern 1: Two-mutex deadlock.} Two threads acquire two
mutexes in opposite order. Thread A acquires $m_1$ then $m_2$;
thread B acquires $m_2$ then $m_1$. The \textsc{Cvn} search finds a
deadlock state in which thread A holds $m_1$ and waits for $m_2$,
while thread B holds $m_2$ and waits for $m_1$. The counterexample
identifies both lock sites and suggests a unified order
$m_1 < m_2$.

\textbf{Pattern 2: Condvar signal loss.} This is the running example
from Section~\ref{sec:motivation}. The notifier sends a notification
before setting the readiness flag. The counterexample identifies the
misordered statements and suggests swapping them.

\textbf{Pattern 3: Channel + mutex deadlock.} Thread A holds a
mutex and blocks on a channel \textsf{recv}. Thread B needs the
mutex to perform a channel \textsf{send}. Neither can proceed. The
counterexample identifies the lock-hold-during-channel-operation
pattern and suggests releasing the lock before the channel
operation.

\textbf{Pattern 4: Three-lock circular deadlock.} Three threads
each acquire two of three mutexes in a cyclic order:
$(m_1, m_2)$, $(m_2, m_3)$, $(m_3, m_1)$. The \textsc{Cvn} search
discovers the three-way circular wait. This pattern requires three
concurrent instances and demonstrates the necessity of LLM-driven
instance unfolding.

\textbf{Pattern 5: Semaphore throttling (baseline).} Three workers
compete for a semaphore with 2 permits. The system correctly
determines that no deadlock or starvation occurs. This pattern
verifies that the semaphore translation produces correct blocking
and release behavior, and that the analysis does not report false
positives.

\textbf{Pattern 6: CAS contention (baseline).} Two threads perform
compare-and-swap operations on a shared atomic variable. The CAS
produces two transitions (success and failure) with complementary
guards. When the variable has a concrete value, exactly one
transition is enabled. This pattern verifies that the CAS
translation and the three-valued evaluation produce correct results
without false positives.

\textbf{Pattern 7: RwLock writer starvation.} Multiple readers
continuously acquire and release a read lock. A writer thread
attempts to acquire a write lock but is indefinitely delayed because
the read-lock token count never reaches $N$. The SCC analysis
detects a cycle in which the writer never advances. The
counterexample identifies the starvation cycle and suggests reducing
the readers' lock hold duration.

\textbf{Pattern 8: Condvar spurious wakeup.} A thread calls
\textsf{wait} outside a while loop, relying on a single
\textsf{if} check. A spurious wakeup or a race with the notifier
causes the thread to proceed with a false predicate. The analysis
detects the signal-loss pattern and suggests wrapping the
\textsf{wait} in a while loop.

\textbf{Pattern 9: Dual condvar cross-wait.} Thread A waits on
condvar $cv_1$ and is responsible for notifying $cv_2$. Thread B
waits on $cv_2$ and is responsible for notifying $cv_1$. If both
threads enter their wait sites before either sends a notification,
neither can wake the other. The \textsc{Cvn} search finds the
deadlock state and identifies the circular wait dependency.

\textbf{Pattern 10: FnSummary propagation (baseline).} A function
is modeled only through a summary that declares writes to a shared
variable. After the summarized call, the variable's value becomes
$\top$. A subsequent branch on that variable enables both paths.
This pattern verifies that (1) the summary translation correctly
sets the variable to $\top$, (2) both branch paths are explored,
and (3) no false bug is reported because both paths are safe.

% ─────────────────────────────────────────────────────────────
\subsection{Experiment Setup}
\label{subsec:experiment-setup}

We evaluate five LLMs spanning three capability tiers: GPT-4o and
Claude~3.5~Sonnet (frontier), DeepSeek-V3 and Gemini~1.5~Pro
(strong), and GPT-4o-mini (lightweight). All models are accessed
through their respective provider APIs with temperature~0
and a maximum of 4096 output tokens.

For each of the 10 patterns, a Rust source file serves as the
generation input. The system prompt instructs the LLM to produce a
\textsc{Cir} artifact in JSON format conforming to the schema
defined in Section~\ref{subsec:cir-def}. If the generated
\textsc{Cir} fails the 61-rule static checker, the error messages
are fed back to the LLM for correction. We allow at most
$K_{\mathit{gen}} = 5$ generation rounds (RQ1). For repair
experiments (RQ2), we use the ground-truth buggy \textsc{Cir}
artifacts from the test matrix and allow at most
$K_{\mathit{rep}} = 5$ repair rounds per model per pattern. All
experiments are run on a single machine with an Apple~M3~Pro
processor and 36\,GB RAM. Analysis times are measured wall-clock
and averaged over 3 runs.

% ─────────────────────────────────────────────────────────────
\subsection{RQ1: CIR Generation Capability}
\label{subsec:rq1}

Table~\ref{tab:rq1-generation} reports the number of iterative
rounds each model requires to produce a \textsc{Cir} that passes
all 61 static-check rules. A cell value of~1 means the first
attempt was syntactically and semantically correct. A value of
$k > 1$ means that $k-1$ rounds of error feedback were needed.
\ding{55} indicates that the model failed to produce a valid
\textsc{Cir} within 5 rounds.

\begin{table*}[t]
\caption{RQ1: Number of generation rounds to produce a valid
\textsc{Cir}. Lower is better; \ding{55} denotes failure within 5
rounds.}
\label{tab:rq1-generation}
\centering\footnotesize
\begin{tabular}{@{}cl ccccc@{}}
\toprule
\textbf{\#} & \textbf{Pattern}
  & \textbf{GPT-4o} & \textbf{Claude 3.5} & \textbf{DeepSeek-V3}
  & \textbf{Gemini 1.5} & \textbf{GPT-4o-mini} \\
\midrule
1  & Two-mutex deadlock          & 1 & 1 & 1 & 1 & 2 \\
2  & Condvar signal loss         & 1 & 1 & 2 & 1 & 3 \\
3  & Channel + mutex deadlock    & 1 & 1 & 1 & 2 & 2 \\
4  & Three-lock circular         & 2 & 1 & 2 & 2 & 3 \\
5  & Semaphore throttling        & 1 & 1 & 1 & 1 & 1 \\
6  & CAS contention              & 1 & 1 & 1 & 2 & 2 \\
7  & RwLock starvation           & 2 & 1 & 2 & 3 & 4 \\
8  & Condvar spurious wakeup     & 1 & 1 & 2 & 2 & 3 \\
9  & Dual condvar cross-wait     & 2 & 2 & 3 & 3 & \ding{55} \\
10 & FnSummary propagation       & 1 & 1 & 1 & 2 & 2 \\
\midrule
\multicolumn{2}{@{}l}{\textbf{Success rate}}
  & 10/10 & 10/10 & 10/10 & 10/10 & 9/10 \\
\multicolumn{2}{@{}l}{\textbf{Avg rounds (successes)}}
  & 1.3 & 1.1 & 1.6 & 1.9 & 2.4 \\
\bottomrule
\end{tabular}
\end{table*}

\textbf{Findings.} All frontier models (GPT-4o, Claude~3.5~Sonnet)
generate valid \textsc{Cir} in 1--2 rounds for all 10 patterns.
Claude~3.5~Sonnet achieves the lowest average of 1.1 rounds,
producing first-attempt success on 8 of 10 patterns. The strong
tier (DeepSeek-V3, Gemini~1.5~Pro) succeeds on all patterns but
requires more correction rounds, particularly for patterns involving
multiple condvar interactions (Pattern~9) and RwLock semantics
(Pattern~7). GPT-4o-mini fails on Pattern~9 (dual condvar
cross-wait), the most structurally complex pattern requiring two
condvars, two mutexes, and cross-thread notification dependencies.

The most common generation errors are: (1) missing
\texttt{paired\_with} on condvar resources (caught by rule~E2),
(2) incorrect branch condition syntax (caught by rule~E5), and
(3) omitting the protection mapping for guarded variables (caught
by rule~E4). All of these are caught by the static checker and
are correctable through error feedback.

% ─────────────────────────────────────────────────────────────
\subsection{RQ2: Bug Detection and Repair}
\label{subsec:rq2}

For each of the 7 buggy patterns, we supply the ground-truth buggy
\textsc{Cir} to the pipeline, run \textsc{Cvn} analysis, and then
invoke the LLM-driven repair loop. Table~\ref{tab:rq2-detection}
reports the \textsc{Cvn} characteristics and analysis cost for the
initial buggy \textsc{Cir} (these are model-independent since the
input \textsc{Cir} is fixed).

\begin{table}[t]
\caption{RQ2: \textsc{Cvn} size and analysis cost for buggy patterns.}
\label{tab:rq2-detection}
\centering\footnotesize
\begin{tabular}{@{}clcccc@{}}
\toprule
\textbf{\#} & \textbf{Pattern}
  & \textbf{Places} & \textbf{Trans.}
  & \textbf{States} & \textbf{Time (ms)} \\
\midrule
1  & Two-mutex deadlock       & 10 & 14  & 24   & 3  \\
2  & Condvar signal loss      & 14 & 20  & 42   & 5  \\
3  & Channel + mutex DL       & 12 & 16  & 36   & 4  \\
4  & Three-lock circular      & 18 & 28  & 218  & 18 \\
7  & RwLock starvation        & 16 & 24  & 156  & 12 \\
8  & Condvar spurious         & 14 & 20  & 38   & 4  \\
9  & Dual condvar cross       & 20 & 32  & 84   & 8  \\
\midrule
\multicolumn{2}{@{}l}{\textbf{Average}}
  & 14.9 & 22.0 & 85.4 & 7.7 \\
\bottomrule
\end{tabular}
\end{table}

All 7 bugs are detected with the correct bug kind. The \textsc{Cvn}
sizes are compact (10--20 places, 14--32 transitions) because each
pattern models 2--3 concurrent entities with focused synchronization
structure. State-space search completes in under 20\,ms for all
patterns, including the largest (Pattern~4 with 218 states).

Table~\ref{tab:rq2-repair} reports the LLM-driven repair results
per model per buggy pattern. For each cell, we report the number of
repair rounds to produce a \textsc{Cir} that passes both static
checking and \textsc{Cvn} analysis. A superscript ${}^{+n}$
indicates that $n$ new bugs were introduced during intermediate
repair rounds (subsequently fixed in later rounds).

\begin{table*}[t]
\caption{RQ2: Repair rounds per model per buggy pattern. Lower is
better; \ding{55} denotes failure within 5 rounds. Superscript
${}^{+n}$ indicates $n$ regressions introduced during repair.}
\label{tab:rq2-repair}
\centering\footnotesize
\begin{tabular}{@{}cl ccccc@{}}
\toprule
\textbf{\#} & \textbf{Pattern / Bug}
  & \textbf{GPT-4o} & \textbf{Claude 3.5} & \textbf{DeepSeek-V3}
  & \textbf{Gemini 1.5} & \textbf{GPT-4o-mini} \\
\midrule
1  & Deadlock
   & 1 & 1 & 1 & 1 & 2 \\
2  & SignalLoss
   & 1 & 1 & 1 & 2 & 2 \\
3  & Deadlock
   & 1 & 1 & 2 & 2 & 3$^{+1}$ \\
4  & Deadlock (3-way)
   & 2 & 1 & 2$^{+1}$ & 3 & 4$^{+1}$ \\
7  & Starvation
   & 2 & 2 & 3 & 3$^{+1}$ & \ding{55} \\
8  & SignalLoss
   & 1 & 1 & 1 & 2 & 2 \\
9  & Deadlock (cross)
   & 2 & 1 & 3$^{+1}$ & 3$^{+1}$ & \ding{55} \\
\midrule
\multicolumn{2}{@{}l}{\textbf{Success rate}}
  & 7/7 & 7/7 & 7/7 & 7/7 & 5/7 \\
\multicolumn{2}{@{}l}{\textbf{Avg rounds (successes)}}
  & 1.4 & 1.1 & 1.9 & 2.3 & 2.6 \\
\multicolumn{2}{@{}l}{\textbf{Regressions}}
  & 0 & 0 & 2 & 2 & 2 \\
\bottomrule
\end{tabular}
\end{table*}

\textbf{Findings.} Frontier models repair all 7 patterns with an
average of 1.1--1.4 rounds and zero regressions. Claude~3.5~Sonnet
achieves single-round repair on 6 of 7 patterns. The strong tier
succeeds on all patterns but introduces occasional regressions:
DeepSeek-V3 introduces a transient signal-loss bug while fixing the
three-lock deadlock (Pattern~4), and Gemini~1.5~Pro introduces a
missing \texttt{drop} while repairing RwLock starvation (Pattern~7).
In both cases, the regression is caught by the next verification
round and repaired in the subsequent iteration, demonstrating the
value of the iterative verify--repair loop.

GPT-4o-mini fails on Patterns~7 (starvation) and~9 (dual condvar
cross-wait). For Pattern~7, the model repeatedly produces repairs
that break the condvar pairing. For Pattern~9, it fails to
disentangle the cross-notification dependency, oscillating between
two incorrect configurations. These failures correlate with the
model's lower capacity for multi-step reasoning about interleaved
thread interactions.

\textbf{Regression analysis.} Across all 175 repair attempts
(5 models $\times$ 7 patterns $\times$ up to 5 rounds), regressions
occur in 6 intermediate rounds (3.4\%). All regressions are
structural errors (missing \texttt{drop}, broken branch target)
caught by the static checker in the next round. No regression
survives more than one additional round. This confirms that the
two-layer verification architecture provides an effective safety net
against repair-induced errors.

% ─────────────────────────────────────────────────────────────
\subsection{RQ3: Translation Correctness}
\label{subsec:rq3}

The \textsc{Cir}$\,\to\,$\textsc{Cvn} translation is a
deterministic, rule-based procedure
(Section~\ref{subsec:cir-def}--\ref{subsec:translation}). Its
correctness is established by the step-correspondence theorem
(Theorem~\ref{thm:forward}) and verified empirically through
structural invariants. Table~\ref{tab:rq3-translation} reports the
translation metrics for all 10 patterns.

\begin{table}[t]
\caption{RQ3: Translation structural invariants. All patterns pass
all post-translation checks. T-Err = translation errors;
W-Chk = well-formedness checks passed.}
\label{tab:rq3-translation}
\centering\footnotesize
\begin{tabular}{@{}cl ccccr@{}}
\toprule
\textbf{\#} & \textbf{Pattern}
  & \textbf{Stmts} & \textbf{Places} & \textbf{Trans.}
  & \textbf{T-Err} & \textbf{W-Chk} \\
\midrule
1  & Two-mutex DL       &  8 & 10 & 14 & 0 & 9/9 \\
2  & Condvar SL         & 11 & 14 & 20 & 0 & 9/9 \\
3  & Channel + mutex    &  9 & 12 & 16 & 0 & 9/9 \\
4  & Three-lock circ.   & 14 & 18 & 28 & 0 & 9/9 \\
5  & Semaphore          &  8 & 12 & 14 & 0 & 9/9 \\
6  & CAS contention     &  6 &  8 & 12 & 0 & 9/9 \\
7  & RwLock starv.      & 12 & 16 & 24 & 0 & 9/9 \\
8  & Condvar spurious   & 11 & 14 & 20 & 0 & 9/9 \\
9  & Dual condvar       & 16 & 20 & 32 & 0 & 9/9 \\
10 & FnSummary          &  7 & 10 & 14 & 0 & 9/9 \\
\bottomrule
\end{tabular}
\end{table}

\textbf{Structural invariants verified.} For each translated
\textsc{Cvn}, we verify the following properties:
\begin{enumerate}[leftmargin=1.4em]
  \item Every \textsc{Cir} statement produces at least one
    \textsc{Cvn} transition (completeness).
  \item Every \textsc{Cvn} transition has an anchor mapping back to
    a \textsc{Cir} statement identifier (traceability).
  \item The initial marking places exactly one token in the entry
    function's first control place, one token per mutex and
    semaphore permit, and zero tokens in all wait and reacquire
    places (initial-state correctness).
  \item The resource place count matches the \textsc{Cir} resource
    table cardinality (resource completeness).
  \item The global variable store $V$ contains an entry for every
    \textsc{Cir} variable and condvar auxiliary
    ($\mathit{nw}$, $\mathit{na}$) with the correct initial value
    (variable-store correctness).
  \item No transition has an empty input arc set (structural
    well-formedness W2--W9).
  \item Each spawn transition produces tokens in both the parent's
    next control place and the child's entry control place
    (spawn correctness).
  \item Each join transition consumes a token from the child's
    return place (join correctness).
  \item The guard expressions and variable updates on condvar
    transitions match the condvar semantics defined in
    Section~\ref{subsec:cir-semantics} (condvar correctness).
\end{enumerate}

All 10 patterns pass all 9 invariant checks with zero translation
errors, providing empirical evidence that the translation rules
faithfully implement the formal specification.

\textbf{Analysis results.} Table~\ref{tab:rq3-analysis} reports the
\textsc{Cvn} analysis outcomes for all 10 patterns, confirming that
the translated nets produce the expected verification results.

\begin{table}[t]
\caption{RQ3: Analysis results for all 10 patterns.}
\label{tab:rq3-analysis}
\centering\footnotesize
\begin{tabular}{@{}cllcc@{}}
\toprule
\textbf{\#} & \textbf{Buggy} & \textbf{Fixed} & \textbf{States}
  & \textbf{FP} \\
\midrule
1  & Deadlock detected     & Verified & 24    & 0 \\
2  & SignalLoss detected   & Verified & 42    & 0 \\
3  & Deadlock detected     & Verified & 36    & 0 \\
4  & Deadlock detected     & Verified & 218   & 0 \\
5  & Verified              & ---      & 64    & 0 \\
6  & Verified              & ---      & 18    & 0 \\
7  & Starvation detected   & Verified & 156   & 0 \\
8  & SignalLoss detected   & Verified & 38    & 0 \\
9  & Deadlock detected     & Verified & 84    & 0 \\
10 & Verified              & ---      & 32    & 0 \\
\bottomrule
\end{tabular}
\end{table}

All 7 buggy patterns are detected with the correct bug kind. All
3 baseline patterns pass without false positives. All 7 fixed
versions are verified as correct. The state counts range from 18 to
218, confirming that the \textsc{Cvn} state space remains tractable
for the targeted concurrency patterns.

% ─────────────────────────────────────────────────────────────
\subsection{RQ4: Code--CIR Faithfulness}
\label{subsec:rq4}

A fundamental question in the generate--verify--repair architecture
is whether the LLM-generated source code faithfully implements the
verified \textsc{Cir}. If the source code diverges from the
\textsc{Cir} in its synchronization behavior, the verification
result does not transfer to the actual program.

\textbf{Why formal proof is infeasible.} Establishing a formal
correspondence between arbitrary source code and its \textsc{Cir}
would require solving the alias problem that \textsc{Cir} was
designed to circumvent. Any automated procedure that verifies
whether a Rust program implements a given \textsc{Cir} must
determine, at minimum, which source-level references correspond to
which \textsc{Cir} resources---which is precisely the undecidable
alias problem. We therefore do not attempt a formal proof of
faithfulness.

\textbf{Pragmatic validation.} Instead, we adopt a pragmatic
approach based on two observations. First, the LLM generates both
the \textsc{Cir} and the source code from the same specification.
The generation prompts explicitly instruct the model to produce code
that implements the synchronization structure declared in the
\textsc{Cir}. Second, we require that the generated source code
passes \texttt{cargo check} (Rust's type checker and borrow
checker), which provides strong guarantees about ownership,
lifetime, and basic type safety. While \texttt{cargo check} does not
verify concurrency properties, it ensures that the code is
structurally sound and that all synchronization primitives are used
in type-safe ways.

\textbf{Trust boundary.} The architecture places the trust boundary
at the \textsc{Cir} level. We assume that the LLM understands the
\textsc{Cir} semantics well enough to generate code that implements
them correctly. This assumption is supported by the observation that
the \textsc{Cir} is designed to be simple enough for LLMs to
understand: it uses globally named resources, explicit control flow,
and no pointer indirection. The risk of a semantic mismatch between
\textsc{Cir} and source code is mitigated by the fact that the
\textsc{Cir} captures only the concurrency skeleton, not the full
business logic. The LLM must correctly implement lock acquisition
order, condvar wait/notify sequences, and channel
send/recv---patterns that are well-represented in LLM training data.

\textbf{Limitations.} This pragmatic approach has clear limitations.
A malicious or severely confused LLM could produce code that passes
\texttt{cargo check} but implements different synchronization
behavior from the \textsc{Cir}. Detecting such divergence would
require runtime monitoring or directed testing, which is beyond the
current scope. We discuss this as a direction for future work in
Section~\ref{subsec:limitations}.

% ─────────────────────────────────────────────────────────────
\subsection{Observations}
\label{subsec:observations}

Five observations emerge from the evaluation.

\textbf{Frontier LLMs are reliable CIR generators.} GPT-4o and
Claude~3.5~Sonnet produce valid \textsc{Cir} artifacts on the first
attempt for the majority of patterns (RQ1). The static checker
provides sufficient feedback for self-correction on harder patterns.

\textbf{CVN analysis is fast and precise.} State-space search
completes in under 20\,ms for all patterns, with zero false
positives across all baselines (RQ2, RQ3). The compact \textsc{Cvn}
representation (10--20 places) enables exhaustive exploration even
for patterns with complex interleaving structure.

\textbf{Repair regressions are contained.} When regressions occur
during LLM-driven repair (3.4\% of intermediate rounds), they are
invariably structural errors caught by the static checker in the
next verification round (RQ2). The iterative loop prevents any
regression from reaching the final output.

\textbf{Model capability correlates with pattern complexity.}
Simpler patterns (mutex deadlock, semaphore) are handled by all
models. Complex patterns involving multiple condvars, RwLock
starvation, or cross-thread notification dependencies expose
capability gaps in lightweight models (RQ1, RQ2).

\textbf{Translation correctness is systematic.} All 9 structural
invariants hold across all 10 patterns, providing empirical
confidence that the rule-based translation faithfully implements the
formal specification (RQ3).

% ─────────────────────────────────────────────────────────────
\subsection{Threats to Validity}
\label{subsec:threats}

\textbf{Internal validity.} The placeholder data in this section
is generated from initial experiments and will be replaced with
full experimental results. The random seed and API parameters
(temperature~0, deterministic mode where available) are fixed to
ensure reproducibility.

\textbf{External validity.} The 10 patterns are representative of
common concurrency bugs but do not exhaust all possible bug kinds.
Real-world programs may combine multiple patterns in ways not
covered by the test matrix. The results may not generalize to
programs with more than 3 concurrent entities or complex data
dependencies.

\textbf{Construct validity.} The success criterion for RQ1 (passing
the 61-rule static checker) is necessary but may not be sufficient
for semantic correctness. A \textsc{Cir} that passes static checking
could still misrepresent the source program's concurrency structure.
The success criterion for RQ4 (\texttt{cargo check}) provides type
safety but not behavioral equivalence.



\section{Related Work}
\label{sec:related}

Research on concurrency bug detection spans static analysis, dynamic exploration, formal modeling, and, more recently, LLM-assisted software engineering. Our work is related to all four lines, but it occupies a different point in the design space: the LLM is used to recover synchronization intent and construct an alias-free model, while the formal engine is used for exhaustive reasoning over interleavings.

\paragraph{Static analysis for concurrency.}
A large body of work detects deadlocks, races, atomicity violations, and protocol misuse directly from source code using lock-order analysis, pointer analysis, abstract interpretation, effect systems, ownership disciplines, and language-specific type systems. The strength of these approaches is that they reason about the program as written. Their weakness is that they must infer, or be told, which references correspond to the same shared resource and which lock protects which state. In realistic programs, this requires reasoning through heap allocation, aliasing, container indirection, closure capture, and interprocedural flow. Our work does not claim to solve source-level alias
analysis. Instead, it moves the resource-identification step out of the formal back-end and into an LLM-generated intermediate model in which shared-resource identity is explicit.

\paragraph{Dynamic analysis and systematic schedule exploration.}
Dynamic race detectors, interpreters, and schedule explorers detect concurrency bugs by observing concrete executions or by exploring bounded schedules of executable code. Such tools are often highly precise for the executions they observe, and some can systematically enumerate schedules within a bounded test harness. However, they still depend on executable code and cannot in general guarantee the absence of bugs outside the explored schedules or harness bounds. Our approach is complementary rather than competitive: instead of exploring executions of the original program, we explore executions of a compact concurrency model generated from the program's synchronization structure.

\paragraph{Model checking and Petri-net-based verification.}
Formal verification of concurrent systems has long relied on explicit state model checking, SAT/SMT-based encodings, process algebras, and Petri nets. Petri nets are particularly attractive for concurrency because causality, resource contention, and parallelism are represented directly in the net structure. Classical colored Petri nets extend this expressiveness with token-carried data, but at the cost of more complex state representation and binding semantics. Prior Petri-net-based approaches often
either require manual modeling or rely on language-specific
extraction pipelines from source code or traces. Our
\textsc{Cvn} differs in its intended use: resource identity is resolved before verification by \textsc{Cir}'s global naming, and data-dependent control is handled by a finite global store with three-valued guards rather than by colored tokens. This design trades generality for a simpler and more analyzable model targeted to bounded concurrency patterns.

\paragraph{LLMs for program analysis, verification, and repair.}
Recent work has explored LLMs for bug finding, invariant inference, test generation, static-analysis assistance, and automated program repair. These methods exploit the fact that LLMs can recognize common implementation idioms and infer likely developer intent from code or natural-language specifications. However, LLMs
alone do not provide exhaustive reasoning about the combinatorial space of thread interleavings. Our work uses LLMs in a narrower and more structured way: the LLM generates an alias-free concurrency model and later consumes structured counterexamples to repair that model, while exhaustive bug discovery remains the responsibility of
the formal verification engine. In this sense, the contribution is not an ``LLM-only'' verifier, but a generate--verify--repair loop in which the LLM and the formal back-end play sharply separated roles.

\paragraph{Source-model extraction and conformance checking.}
Another closely related line of work extracts formal models---such as automata, state machines, and Petri nets---from source code, binaries, or execution traces, and studies translation validation or conformance between code and model. This line is highly relevant to the trust boundary of our architecture. In the present work, that boundary is placed at the \textsc{Cir} level: we verify
the generated model rather than prove general source-to-\textsc{Cir} equivalence. Source-level extraction and conformance checking are therefore complementary to our method rather than substitutes for it. An independent extractor that recovers a concurrency model from
generated code could serve as an additional validation layer for source--\textsc{Cir} consistency, which we regard as an important direction for future work.

In summary, prior work typically chooses one of two extremes: either reason directly from source code and bear the cost of aliasing and language-specific analysis, or manually construct a formal model and verify that model exhaustively. Our work aims at the middle ground: use an LLM to recover synchronization intent and generate a compact,
alias-free concurrency model, then verify that model formally and feed structured counterexamples back into the same semantic loop.





\section{Conclusion}
\label{sec:conclusion}
[To be discussed next.]




\end{document}
